\chapter{Real blow-up and manifolds with corners}\label{ch:blowup}
\providecommand{\Vb}{\mathcal{V}_{\mathrm{b}}}
\providecommand{\bT}{{}^{\mathrm{b}}T}
\providecommand{\relRC}[2]{{}^{#1}\overline{#2}}
\providecommand{\oRC}[1]{{}^{1}\overline{#1}}

Several spaces in this book are built by ``introducing polar coordinates'' around a
submanifold: the radial compactification of a vector space in \cref{ch:distributions},
the blow-up of the origin used for homogeneous distributions in \cref{ch:conormal}, the
semiclassical and adiabatic double spaces of \cref{ch:semiclassical}, the relative
compactifications carrying product-type symbols in \cref{ch:product}, and the spaces on
which the symbols of \cref{ch:families,ch:eta} live. This appendix gives the general
construction, called \emph{real} (or \emph{radial}) \emph{blow-up}, and proves the facts
about it that those chapters use.

Blow-up is useful for three closely related purposes. It \emph{resolves singular
functions}: $\abs z$ is not smooth on $\R^n$, but in polar coordinates it is the
coordinate $r$. It \emph{resolves singular spaces}: the cone $\{t^2=x^2+y^2,\ t\ge0\}$
becomes, in polar coordinates around the origin of $\R^3$, the smooth manifold with
boundary $[0,\infty)\times\Sphere^1$. And it \emph{resolves degenerate vector fields}:
the linear vector fields $z_i\pa_{z_j}$ vanish at $0$, but they lift to polar coordinates
as smooth vector fields which, together, span all vector fields tangent to the new
boundary. The price is that each blow-up creates a boundary hypersurface, so after two
blow-ups we meet corners. We therefore begin with manifolds with corners
(\cref{blowup:sec:corners}), blow up points (\cref{blowup:sec:point}) and then general
product-type submanifolds (\cref{blowup:sec:psub}). We then study functions and
densities on blown-up spaces, including the resolution of semiclassical kernels
(\cref{blowup:sec:functions}), the behaviour of submanifolds under blow-up and the
commutation of blow-ups (\cref{blowup:sec:commute}), projective blow-up and the blow-up
of hypersurfaces (\cref{blowup:sec:projective}), and the maps which replace fibrations
in this category (\cref{blowup:sec:bfib}). \Cref{blowup:sec:dictionary} lists, with
references, where each of the spaces used in the book comes from.

\begin{discussion}
Why not simply ``use polar coordinates''? Because polar coordinates depend on a choice
of linear coordinates, and we need an object which is defined for a submanifold of a
manifold, with no choices. The central fact of the theory
(\cref{blowup:thm:invariance}) is that every diffeomorphism preserving the centre of
the blow-up lifts to a diffeomorphism of the polar-coordinate spaces. Everything else
---the global construction, its independence of choices, the lifting of vector
fields--- follows from this one lemma and from explicit computations in polar or
projective coordinates.
\end{discussion}

%======================================================================
\section{Manifolds with corners}\label{blowup:sec:corners}

Manifolds with boundary were introduced in \cref{ch:boundary}
(\cref{boundary:sec:manifolds}). We now allow several boundary hypersurfaces to meet,
transversally, at corners.

\subsection{Local models and smooth maps}
For $0\le k\le n$ the model space is
\begin{equation*}
  \R^{n,k}=[0,\infty)^k\times\R^{n-k},
\end{equation*}
with coordinates $(x_1,\dots,x_k,y_1,\dots,y_{n-k})$; the $x_i$ are \emph{boundary
coordinates}. As in \cref{boundary:sec:manifolds}, an open set in $\R^{n,k}$ means a
relatively open set $O=O'\cap\R^{n,k}$, $O'\subset\R^n$ open, and a map or function on
$O$ is \emph{smooth} if it is the restriction of a smooth map or function on such an
$O'$. A diffeomorphism is a smooth bijection with smooth inverse.

\begin{remark}\label{blowup:rem:smooth}
As for manifolds with boundary (\cref{boundary:prop:extension}), $u$ is smooth on $O$ if
and only if it is smooth on $O\cap(0,\infty)^k\times\R^{n-k}$ and all its derivatives
are bounded on $K\cap(0,\infty)^k\times\R^{n-k}$ for each compact $K\subset O$. The
proof of \cref{boundary:prop:extension} applies to one boundary coordinate at a time:
Borel's lemma with parameters extends $u$ across $\{x_1=0\}$, with the same bounds on
compact sets, and then one continues with $x_2$, and so on. We use this criterion only
in \cref{blowup:lem:flat}.
\end{remark}

For a point $p\in\R^{n,k}$ let $L(p)\subset\{1,\dots,k\}$ be the set of $i$ with
$x_i(p)=0$; its size $|L(p)|$ is the \emph{codimension} of $p$. The \emph{inward cone}
$C_p$ is the set of velocities $\gamma'(0)$ of smooth curves
$\gamma:[0,\epsilon)\to\R^{n,k}$ with $\gamma(0)=p$. Clearly
\begin{equation*}
  C_p=\{v\in\R^n;\ v_i\ge0\ \text{for}\ i\in L(p)\}.
\end{equation*}

\begin{lemma}\label{blowup:lem:codim}
A diffeomorphism $F:O_1\to O_2$ between open subsets of $\R^{n,k}$ and $\R^{n,k'}$
satisfies $DF(p)(C_p)=C_{F(p)}$, and preserves the codimension of points.
\end{lemma}
\begin{proof}
$F$ maps curves starting at $p$ to curves starting at $F(p)$, and $F^{-1}$ does the
converse; by the chain rule (for smooth extensions) $DF(p)C_p=C_{F(p)}$. The largest
linear subspace contained in $C_p$ is $\{v;\ v_i=0,\ i\in L(p)\}$, of dimension
$n-|L(p)|$. The linear isomorphism $DF(p)$ maps it onto the corresponding subspace of
$C_{F(p)}$, so $|L(p)|=|L(F(p))|$.
\end{proof}

\subsection{Manifolds with corners}
A \emph{chart with corners} on a topological space $X$ is a homeomorphism of an open
set $U\subset X$ onto an open subset of some $\R^{n,k}$ (with $n$ fixed and $k$ allowed
to vary). A $t$-manifold (``manifold with corners in the weak sense'') is a
second-countable Hausdorff space with a maximal atlas of such charts whose transition
maps are smooth. By \cref{blowup:lem:codim} the codimension of a point is well
defined. The points of codimension $0$ form the interior $X^\circ$, and the points of
codimension $1$ form a manifold $X^{(1)}$ without boundary of dimension $n-1$. The
\emph{boundary hypersurfaces} of $X$ are the closures of the components of $X^{(1)}$.
The tangent bundle, vector fields, densities and differential operators are defined
exactly as for manifolds without boundary, and $C_pX\subset T_pX$ denotes the inward
cone.

\begin{definition}\label{blowup:def:mwc}
A \emph{manifold with corners} is a $t$-manifold $X$ each of whose boundary
hypersurfaces $H$ is \emph{embedded}: there is a \emph{boundary defining function}
$\rho_H\in\CI(X)$ with $\rho_H\ge0$, $H=\{\rho_H=0\}$ and $d\rho_H\ne0$ at every point
of $H$. We write $\mathcal M_1(X)$ for the set of boundary hypersurfaces.
\end{definition}

\begin{figure}[ht]
\centering
\begin{tikzpicture}[>=Stealth,scale=0.95]
  \begin{scope}
    \fill[gray!15] (0,0) rectangle (2.4,2.4);
    \draw[very thick] (0,2.4)--(0,0)--(2.4,0);
    \fill (0,0) circle (1.8pt);
    \node[left,font=\small] at (0,1.3) {$H_1=\{x_1=0\}$};
    \node[below,font=\small] at (1.4,0) {$H_2=\{x_2=0\}$};
    \node[font=\small] at (1.4,1.4) {$\R^{2,2}$};
    \node[font=\small] at (1.2,-1.0) {(a) a corner};
  \end{scope}
  \begin{scope}[xshift=4.6cm]
    \fill[gray!15] (0,0) .. controls (2.2,0.2) and (2.6,2.6) .. (1.2,2.5)
       .. controls (0.2,2.4) and (0.2,2.2) .. (0,0);
    \draw[very thick] (0,0) .. controls (2.2,0.2) and (2.6,2.6) .. (1.2,2.5)
       .. controls (0.2,2.4) and (0.2,2.2) .. (0,0);
    \fill (0,0) circle (1.8pt);
    \node[font=\small,align=center] at (1.2,-1.0) {(b) teardrop: one\\ hypersurface, not embedded};
  \end{scope}
  \begin{scope}[xshift=8.8cm]
    \fill[gray!15] (0,0) rectangle (2.4,2.4);
    \draw[very thick] (0,0) rectangle (2.4,2.4);
    \draw[thick,dashed] (0,0)--(2.4,2.4);
    \node[font=\small] at (1.75,0.9) {$\Diag$};
    \fill (0,0) circle (1.8pt);
    \node[font=\small,align=center] at (1.2,-1.0) {(c) $\Diag\subset[0,1]^2$ is\\ not a p-submanifold};
  \end{scope}
\end{tikzpicture}
\caption{(a) Near a corner of codimension two, the two boundary hypersurfaces are
coordinate hyperplanes. (b) The teardrop is a $t$-manifold but not a manifold with
corners: its only boundary hypersurface meets itself at the corner, where no defining
function can have non-zero differential. (c) The diagonal of the square is an embedded
submanifold but not of product type at $(0,0)$; blowing up the corner repairs this
(\cref{blowup:fig:corner}).}
\label{blowup:fig:corners}
\end{figure}

\begin{lemma}\label{blowup:lem:bdf}
Let $X$ be a manifold with corners and $p\in X$ of codimension $l$, with chart
$(x,y)\in\R^{n,l}$ centred at $p$. Then exactly $l$ boundary hypersurfaces
$H_1,\dots,H_l$ pass through $p$, they can be numbered so that $H_i=\{x_i=0\}$ near
$p$, and $\rho_{H_i}=x_ig_i$ with $g_i$ smooth and positive near $p$. Any two boundary
defining functions of $H$ differ by a positive smooth factor.
\end{lemma}
\begin{proof}
Near $p$ the codimension-one points form the sets $\{x_i=0,\ x_j>0\ (j\ne i)\}$, so each
local hypersurface $\{x_i=0\}$ lies in some boundary hypersurface. If $\{x_1=0\}$ and
$\{x_2=0\}$ lay in the same $H$, then $\rho_H$ would vanish on both, hence
$d\rho_H(p)=0$, which is excluded. Taylor's formula gives $\rho_{H_i}=x_ig_i$ with
$g_i=\int_0^1\pa_{x_i}\rho_{H_i}(\dots,tx_i,\dots)\,dt$. On $\{x_i=0\}$ we have
$d\rho_{H_i}=g_i\,dx_i\neq0$, so $g_i\ne0$ there; for $x_i>0$, $g_i=\rho_{H_i}/x_i>0$.
Hence $g_i>0$ near $p$. The quotient of two defining functions is therefore smooth and
positive near $H$, and it is smooth and positive away from $H$.
\end{proof}

The intersection of any collection of boundary hypersurfaces is, near each of its
points, of the form $\{x_i=0,\ i\in I\}$; its components are the \emph{boundary faces}
of $X$, and are manifolds with corners. The product of two manifolds with corners is a
manifold with corners, with boundary hypersurfaces $H_1\times X_2$ and $X_1\times H_2$.
Examples: $[0,1]^n$, the closed half-sphere $\RC{\R^n}$ of \cref{ch:distributions},
and the \emph{corner spheres}
\begin{equation*}
  \Sphere^{a,b}=\{\omega\in\R^{a+b};\ \abs\omega=1,\ \omega_1,\dots,\omega_a\ge0\},
\end{equation*}
which are the unit spheres of the models $\R^{a+b,a}$. Thus $\Sphere^{0,k}=\Sphere^{k-1}$,
$\Sphere^{1,k}$ is a closed half-sphere and $\Sphere^{2,0}$ is a quarter circle. Near a
point $\omega\in\Sphere^{a,b}$ choose $m$ with $\omega_m\ne0$; then the
$\omega_k$, $k\ne m$, are coordinates, with $\omega_k\ge0$ for $k\le a$, $k\neq m$, so
$\Sphere^{a,b}$ is a manifold with corners whose boundary hypersurfaces
$\{\omega_i=0\}$, $i\le a$, have the defining functions $\omega_i$.

\subsection{Product-type submanifolds}
Embedded submanifolds of manifolds with corners can meet the boundary in complicated
ways. The ones we can blow up are those that look like coordinate subspaces.

\begin{definition}\label{blowup:def:psub}
A closed subset $Y$ of a manifold with corners $X$ is a \emph{p-submanifold}
(product-type submanifold) if each $p\in Y$ has a chart of the form
\begin{equation}\label{blowup:eq:adapted}
  (x,y)\in\R^{a+b,a}\times\R^{m,c},\qquad Y=\{x=0\}\ \text{near}\ p.
\end{equation}
Such a chart is \emph{adapted} to $Y$; $x$ are \emph{normal coordinates}, and $(a,b)$ is
the \emph{normal type} of $Y$ at $p$ (it is locally constant). $Y$ is an \emph{interior}
p-submanifold if $a=0$ everywhere. Two p-submanifolds $Y$, $Z$ are \emph{joint} if each
point of $Y\cap Z$ has a chart $(z_1,\dots,z_n)\in\R^{n,k}$ in which
$Y=\{z_i=0,\ i\in A\}$ and $Z=\{z_i=0,\ i\in B\}$ for index sets $A$, $B$; they are
\emph{transversal} if moreover $A\cap B=\emptyset$ in such charts.
\end{definition}

A p-submanifold is a manifold with corners, with coordinates $y$; its boundary
hypersurfaces are the components of the sets $Y\cap H$, $H\in\mathcal M_1(X)$, $Y\not\subset
H$, with defining functions $\rho_H|_Y$. Every boundary face $F$ of codimension $a$ is a
p-submanifold of normal type $(a,0)$. In an adapted chart the boundary hypersurfaces
through $p$ are the coordinate hyperplanes of boundary coordinates; those containing
$Y$ are $\{x_i=0\}$, $i\le a$.

The \emph{normal bundle} of $Y$ is $NY=T_YX/TY$; in an adapted chart, $N_pY\cong\R^{a+b}$
by the components $dx(v)$ of $v\in T_pX$. The \emph{inward normal cone} $N^+_pY$ is the
image of $C_pX$ in $N_pY$; in an adapted chart it is $[0,\infty)^a\times\R^b$. Finally
\begin{equation*}
  S^+_pY=(N^+_pY\setminus0)/\R^+\cong\Sphere^{a,b},\qquad S^+NY=\bigcup_{p\in Y}S^+_pY,
\end{equation*}
the \emph{inward-pointing spherical normal bundle}. For interior $Y$ this is the full
sphere bundle $SNY=(NY\setminus0)/\R^+$.

\begin{example}
The diagonal of $[0,1]^2$ is an embedded submanifold but not a p-submanifold
(\cref{blowup:fig:corners}(c), \cref{blowup:ex:diag-not-p}). The curve $\{x_2=x_1^2\}$ in
$\R^{2,1}=[0,\infty)\times\R$ is not a p-submanifold at $0$ either: it is tangent to the
boundary there.
\end{example}

\subsection{b-vector fields and b-maps}
\begin{lemma}\label{blowup:lem:Vb}
For a smooth vector field $V$ on a manifold with corners $X$ the following are
equivalent: \textup{(i)} $V$ is tangent to every boundary hypersurface; \textup{(ii)}
$V\rho_H\in\rho_H\CI(X)$ for every $H\in\mathcal M_1(X)$; \textup{(iii)} in every chart,
$V=\sum_ia_ix_i\pa_{x_i}+\sum_jb_j\pa_{y_j}$ with smooth $a_i,b_j$. The space $\Vb(X)$ of
such vector fields is a Lie algebra and a $\CI(X)$-module.
\end{lemma}
\begin{proof}
Write $V=\sum c_i\pa_{x_i}+\sum b_j\pa_{y_j}$. Tangency to $\{x_i=0\}$ means $c_i=0$ there,
i.e.\ $c_i=x_ia_i$ by Taylor's formula; this gives (i)$\iff$(iii), and (ii) follows from
(iii) and \cref{blowup:lem:bdf}, while (ii) implies $c_i=Vx_i\in x_i\CI$ near $p$. If $V,W$
satisfy (ii) then $[V,W]\rho_H=V(W\rho_H)-W(V\rho_H)\in\rho_H\CI$, since
$V(\rho_Ha)=\rho_HVa+aV\rho_H$.
\end{proof}

\begin{definition}\label{blowup:def:bmap}
A smooth map $f:X\to X'$ between manifolds with corners is an \emph{interior b-map} if
for each $H'\in\mathcal M_1(X')$
\begin{equation}\label{blowup:eq:bmap}
  f^*\rho_{H'}=a_{H'}\prod_{H\in\mathcal M_1(X)}\rho_H^{e(H,H')},\qquad
  0<a_{H'}\in\CI(X),\ e(H,H')\in\N_0.
\end{equation}
The integers $e(H,H')$ form the \emph{exponent matrix}. A \emph{b-map} is a map which,
for each $H'$, either satisfies \eqref{blowup:eq:bmap} or has $f^*\rho_{H'}\equiv0$.
\end{definition}

By \cref{blowup:lem:bdf} the condition does not depend on the choice of defining
functions. An interior b-map maps $X^\circ$ into $X'^\circ$, and $f(H)\subset H'$ when
$e(H,H')>0$. Composites of interior b-maps are interior b-maps: substitute
\eqref{blowup:eq:bmap} into itself. Diffeomorphisms, projections
$X_1\times X_2\to X_1$, and (\cref{blowup:prop:props}) blow-down maps are interior
b-maps. The inclusion of a boundary face is a b-map which is not interior.

%======================================================================
\section{Blow-up of a point}\label{blowup:sec:point}

\subsection{Polar coordinates}
The basic example of a blow-up is the polar coordinate map
\begin{equation*}
  \beta:[0,\infty)\times\Sphere^{n-1}\ni(r,\omega)\longmapsto r\omega\in\R^n.
\end{equation*}
It is smooth and surjective, a diffeomorphism of $(0,\infty)\times\Sphere^{n-1}$ onto
$\R^n\setminus0$, and it collapses the \emph{front face} $\{r=0\}$ to the origin
(\cref{blowup:fig:point}). We write $[\R^n;\{0\}]=[0,\infty)\times\Sphere^{n-1}$ for the
moment; \cref{blowup:sec:psub} makes this invariant.

\begin{figure}[ht]
\centering
\begin{tikzpicture}[>=Stealth,scale=0.95]
  \begin{scope}
    \fill[gray!12] (0,0) circle (1.8);
    \foreach \a/\s in {20/solid,75/dashed,135/densely dotted}{
      \draw[thick,\s] (\a:1.8)--(\a+180:1.8);}
    \fill (0,0) circle (2pt);
    \node[font=\small] at (0.35,-0.25) {$0$};
    \node[font=\small] at (0,-2.3) {$\R^2$};
  \end{scope}
  \draw[<-] (2.3,0)--(3.5,0) node[midway,above,font=\small]{$\beta$};
  \begin{scope}[xshift=6cm]
    \fill[gray!12] (0,0) circle (1.8);
    \fill[white] (0,0) circle (0.7);
    \draw[very thick] (0,0) circle (0.7);
    \foreach \a/\s in {20/solid,75/dashed,135/densely dotted}{
      \draw[thick,\s] (\a:0.7)--(\a:1.8);
      \draw[thick,\s] (\a+180:0.7)--(\a+180:1.8);
      \fill (\a:0.7) circle (1.5pt); \fill (\a+180:0.7) circle (1.5pt);}
    \node[font=\small,align=center] at (0,0) {ff};
    \node[font=\small] at (0,-2.3) {$[\R^2;\{0\}]=[0,\infty)\times\Sphere^1$};
  \end{scope}
\end{tikzpicture}
\caption{Blowing up the origin of $\R^2$. Lines through $0$ (left) all meet at the
origin; their lifts (right) meet the front face $\mathrm{ff}=\{r=0\}$ at distinct points,
one for each direction. The lifts are disjoint p-submanifolds.}
\label{blowup:fig:point}
\end{figure}

\begin{lemma}\label{blowup:lem:vf-point}
Every smooth vector field on $\R^n$ which vanishes at $0$ is $\beta$-related to a unique
smooth vector field on $[0,\infty)\times\Sphere^{n-1}$, tangent to the front face. On
$(0,\infty)\times\Sphere^{n-1}$,
\begin{equation}\label{blowup:eq:lin-lift}
  z_i\pa_{z_j}=\omega_i\omega_j\,r\pa_r+\omega_iW_j,\qquad
  W_j=\sum_k(\delta_{jk}-\omega_j\omega_k)\pa_{\omega_k},
\end{equation}
where $W_j$ is the orthogonal projection of $e_j$ onto $T_\omega\Sphere^{n-1}$, a
smooth vector field on the sphere. The lifts span $\Vb([0,\infty)\times\Sphere^{n-1})$
over $\CI([0,\infty)\times\Sphere^{n-1})$. On the other hand
$\pa_{z_j}=\omega_j\pa_r+r^{-1}W_j$ does not lift smoothly.
\end{lemma}
\begin{proof}
For $r>0$, $z_i\pa_{z_j}r=z_iz_j/r=r\omega_i\omega_j$ and
$z_i\pa_{z_j}\omega_k=z_i(\delta_{jk}r^{-1}-z_jz_kr^{-3})=\omega_i(\delta_{jk}-\omega_j\omega_k)$,
which gives \eqref{blowup:eq:lin-lift}; it is smooth up to $r=0$ and tangent to it. A
vector field $V=\sum_jc_j\pa_{z_j}$ vanishing at $0$ has, by Taylor's formula,
$c_j=\sum_iz_ic_{ij}$ with $c_{ij}$ smooth, so $V=\sum c_{ij}z_i\pa_{z_j}$, and its lift
is $\sum(\beta^*c_{ij})\,\omega_i\omega_jr\pa_r+\omega_iW_j$, which is smooth. The lift
is unique because $\beta$ is a diffeomorphism on the dense set $r>0$. Summing
\eqref{blowup:eq:lin-lift} over $i=j$ gives $\sum_iz_i\pa_{z_i}=r\pa_r$, and the
rotations lift to $z_i\pa_{z_j}-z_j\pa_{z_i}=\omega_iW_j-\omega_jW_i$, whose value at
$\omega$ is $\omega_ie_j-\omega_je_i$. For $v\perp\omega$,
$\sum_{i,j}\omega_iv_j(\omega_ie_j-\omega_je_i)=v-(\omega\cdot v)\omega=v$, so the rotations
span $T_\omega\Sphere^{n-1}$ at every point. Since $\Vb([0,\infty)\times\Sphere^{n-1})$
is spanned locally by $r\pa_r$ and vector fields on the sphere, the lifts span it.
Finally $\pa_{z_j}=\omega_j\pa_r+r^{-1}W_j$ by the same computation, and $W_j\ne0$ at
points with $\omega_j\ne\pm1$.
\end{proof}

\subsection{Invariance}
We now prove that polar coordinates do not depend on the choice of linear coordinates,
or even on the choice of smooth coordinates, and we do this in the generality needed
later: with parameters and at corners. Let $\R^{a+b,a}\ni x$ be the normal model and
$P=\R^{m,c}\ni y$ a model for parameters. Define
\begin{equation*}
  \beta:[0,\infty)\times\Sphere^{a,b}\times P\ni(r,\omega,y)\longmapsto(r\omega,y)\in
  \R^{a+b,a}\times P.
\end{equation*}
For $O\subset\R^{a+b,a}\times P$ open write $\tilde O=\beta^{-1}(O)$.

\begin{theorem}\label{blowup:thm:invariance}
Let $O,O'\subset\R^{a+b,a}\times P$ be open and $F:O\to O'$ a diffeomorphism with
$F(O\cap\{x=0\})=O'\cap\{x=0\}$. Write $F=(F_N,F_P)$ and
$A(y)=D_xF_N(0,y)$. Then there is a unique diffeomorphism $\tilde F:\tilde O\to\tilde O'$
with $\beta\circ\tilde F=F\circ\beta$, and on the front face
\begin{equation}\label{blowup:eq:ff-action}
  \tilde F(0,\omega,y)=\Big(0,\frac{A(y)\omega}{\abs{A(y)\omega}},F_P(0,y)\Big).
\end{equation}
\end{theorem}
\begin{proof}
Uniqueness holds because $\beta$ is a diffeomorphism of the dense set $\{r>0\}$ onto
$O\setminus\{x=0\}$; there $\tilde F=\beta^{-1}\circ F\circ\beta$. It remains to show that
this map extends smoothly to $r=0$, with \eqref{blowup:eq:ff-action}; the same argument
for $F^{-1}$ gives a smooth inverse, since the composites are the identity on the dense
set $r>0$.

Fix $(0,y_0)\in O$. Choose a smooth extension $F'$ of $F$ to a neighbourhood
$\{\abs x<\delta\}\times B$ of $(0,y_0)$ in $\R^{a+b}\times\R^m$, with $B$ a ball. The
derivatives of $F$ at points of $O$ are those of $F'$, and since the extension of $F^{-1}$
composed with $F'$ is the identity on an open subset of the model with non-empty
interior near $(0,y_0)$, $DF'(0,y_0)$ is invertible. As $F_N(0,y)=0$ for
$y\in P$ near $y_0$, and $P$ has interior points near $y_0$, $D_yF_N(0,y_0)=0$; so
$DF'(0,y_0)$ is block triangular and $A(y_0)=D_xF_N(0,y_0)$ is invertible. Taylor's
formula in $x$ gives
\begin{equation*}
  F_N'(x,y)=F_N'(0,y)+A(x,y)x,\qquad A(x,y)=\int_0^1D_xF_N'(tx,y)\,dt,
\end{equation*}
and $F_N'(0,y)=0$ for $y\in P$. Shrinking $\delta$ and $B$ we may assume $A(x,y)$ is
invertible. Define, for $\abs r<\delta$, $\omega\in\Sphere^{a+b-1}$, $y\in B$,
\begin{equation*}
  \tilde F'(r,\omega,y)=\Big(r\abs{A(r\omega,y)\omega},\
  \frac{A(r\omega,y)\omega}{\abs{A(r\omega,y)\omega}},\ F_P'(r\omega,y)\Big),
\end{equation*}
which is smooth. For $r>0$ and $(r\omega,y)\in O$ we have
\begin{equation*}
  \beta(\tilde F'(r,\omega,y))=(rA(r\omega,y)\omega,F_P(r\omega,y))=F(r\omega,y),
\end{equation*}
so
$\tilde F'=\beta^{-1}F\beta$ there. In particular $\tilde F'$ maps
$\tilde O\cap\{r>0\}$ into the model $[0,\infty)\times\Sphere^{a,b}\times P$, which is
closed, so by continuity it maps $\tilde O$ near $\{0\}\times\Sphere^{a,b}\times\{y_0\}$
into it. Thus $\beta^{-1}F\beta$ extends smoothly near every point of $\{r=0\}$, and at
$r=0$ the formula reduces to \eqref{blowup:eq:ff-action}.
\end{proof}

\begin{discussion}
The proof uses only that $F$ maps the centre $\{x=0\}$ to itself. Taylor's formula then
says that, to first order, $F$ acts on the normal variables by the invertible matrix
$A(x,y)$, and a linear map acts smoothly on directions. Note that the action on the
front face is through the \emph{normal differential} $A(y)$; this is what makes the
front face canonically the spherical normal bundle in \cref{blowup:thm:blowup}. With
$a=b=0$ in the parameter-free case, and $F$ a change of coordinates fixing $0$, the
theorem is the statement of \cref{distributions:ex:polar}.
\end{discussion}

\subsection{Projective coordinates}
Polar coordinates are awkward for computation, since $\omega$ is constrained to the
sphere. Projective coordinates are often more convenient.

\begin{lemma}\label{blowup:lem:projective}
Let $1\le j\le a+b$ and $\sigma\in\{\pm1\}$, with $\sigma=+1$ if $j\le a$. On the open set
$\{\sigma\omega_j>0\}\subset[0,\infty)\times\Sphere^{a,b}$ the functions
\begin{equation}\label{blowup:eq:projective}
  s=\abs{x_j}=\sigma r\omega_j,\qquad t_k=\frac{x_k}{\abs{x_j}}=\frac{\omega_k}{\sigma\omega_j}
  \quad(k\neq j)
\end{equation}
are coordinates, identifying $\{\sigma\omega_j>0\}$ with
$\{(s,t);\ s\ge0,\ t_k\ge0\ \text{for}\ k\le a\}$. The front face is $\{s=0\}$, and the
lifted boundary hypersurfaces $\{\omega_k=0\}$, $k\le a$, $k\ne j$, are $\{t_k=0\}$. These
sets, for all admissible $(j,\sigma)$, cover the front face.
\end{lemma}
\begin{proof}
The inverse is $\omega=\theta/\abs\theta$, $r=s\abs\theta$, where $\theta$ has $\sigma$ in
slot $j$ and $t_k$ in slot $k$; both maps are smooth. Every $\omega\in\Sphere^{a,b}$ has
some $\omega_j\ne0$, and if $j\le a$ then $\omega_j>0$.
\end{proof}

In projective coordinates, $x_j=\sigma s$ and $x_k=st_k$. Hence, with $\sigma=+1$ for
simplicity, $\pa_{x_k}=s^{-1}\pa_{t_k}$ for $k\ne j$ and
$\pa_{x_j}=\pa_s-s^{-1}\sum_{r\ne j}t_r\pa_{t_r}$, and the linear vector fields lift to
\begin{equation}\label{blowup:eq:proj-lifts}
  x_l\pa_{x_m}\longmapsto
  \begin{cases}
    t_l\pa_{t_m}, & l\ne j,\ m\ne j,\\
    \pa_{t_m}, & l=j,\ m\ne j,\\
    t_l\big(s\pa_s-\sum_{r\ne j}t_r\pa_{t_r}\big), & l\ne j,\ m=j,\\
    s\pa_s-\sum_{r\ne j}t_r\pa_{t_r}, & l=m=j.
  \end{cases}
\end{equation}
For $\sigma=-1$ the second and third lines change sign. These formulae make the
smoothness of the lifts obvious, and they are the main computational tool below. In
particular the Euler field $\sum_lx_l\pa_{x_l}$ lifts to $s\pa_s$.

%======================================================================
\section{Blow-up of p-submanifolds}\label{blowup:sec:psub}

\subsection{The construction}
Let $Y$ be a closed p-submanifold of a manifold with corners $X$. As a set, the
blow-up of $X$ along $Y$ is
\begin{equation*}
  [X;Y]=(X\setminus Y)\sqcup S^+NY,
\end{equation*}
with the \emph{blow-down map} $\beta:[X;Y]\to X$ equal to the identity on $X\setminus Y$
and to the bundle projection $S^+NY\to Y$ on $S^+NY$. The \emph{front face} is
$\mathrm{ff}=\beta^{-1}(Y)=S^+NY$. For an adapted chart $\phi=(x,y):U\to O$ as in
\eqref{blowup:eq:adapted} define the \emph{polar chart}
\begin{gather*}
  \tilde\phi:\beta^{-1}(U)\to\tilde O\subset[0,\infty)\times\Sphere^{a,b}\times\R^{m,c},
  \\
  \tilde\phi(q)=\begin{cases}\big(\abs x,\ x/\abs x,\ y\big)(q),& q\in U\setminus Y,\\
  \big(0,\ \xi/\abs\xi,\ y(p)\big),&q=[v]\in S^+_pY,\ \xi=dx_p(v).\end{cases}
\end{gather*}
It is a bijection, and $\beta$ becomes $(r,\omega,y)\mapsto(r\omega,y)$.

\begin{theorem}\label{blowup:thm:blowup}
There is a unique structure of manifold with corners on $[X;Y]$ for which the charts of
$X\setminus Y$ and all polar charts are charts. With it, $\beta$ is smooth and proper, it
restricts to a diffeomorphism $[X;Y]\setminus\mathrm{ff}\to X\setminus Y$, and the front face
is a union of boundary hypersurfaces (usually one), which is the total space of a fibre bundle over $Y$ with fibre $\Sphere^{a,b}$.
\end{theorem}
\begin{proof}
Transition maps between charts of $X\setminus Y$ are those of $X$. Between a polar chart
and a chart of $X\setminus Y$ the transition is a restriction of $\beta$ or its inverse to
$r>0$, a diffeomorphism. Between two polar charts $\tilde\phi_1$, $\tilde\phi_2$ the
transition is $\widetilde F$ for $F=\phi_2\circ\phi_1^{-1}$, by
\cref{blowup:thm:invariance}. On $r>0$ this is clear. On $r=0$, for $v\in T_pX$ the chain
rule gives $dx^{(2)}_p(v)=A(y)\,dx^{(1)}_p(v)$, because $D_yF_N=0$ at $x=0$; so the
transition acts on $S^+_pY$ by \eqref{blowup:eq:ff-action}, consistently with the
definition of the polar charts. So we
have a smooth atlas; we give $[X;Y]$ the topology for which the charts are homeomorphisms
onto open sets. It is Hausdorff: points of $X\setminus Y$ and of $S^+NY$ with distinct
images in $X$ are separated by preimages of disjoint open sets of $X$, and points of
$S^+_pY$ by a polar chart. It is second countable since countably many charts cover it.
In the charts $\beta$ is $(r,\omega,y)\mapsto(r\omega,y)$ or the identity, hence smooth,
and $\beta^{-1}$ of a compact set is compact (it is so in each polar chart since
$\Sphere^{a,b}$ is compact). In polar charts $\mathrm{ff}=\{r=0\}$ and $r$ is a boundary
coordinate, so the front face is a union of boundary hypersurfaces, and the charts
exhibit it as a bundle over $Y$ with fibre $\Sphere^{a,b}$.

It remains to check that boundary hypersurfaces are embedded, which we do in
\cref{blowup:prop:props} by writing down defining functions. Uniqueness of the
structure is clear, since the charts are prescribed.
\end{proof}

\begin{proposition}\label{blowup:prop:props}
Let $Y\subset X$ be a closed p-submanifold and $\beta:[X;Y]\to X$ the blow-down map.
\begin{enumerate}[label=\textup{(\roman*)}]
\item There is $q_Y\in\CI(X)$, $q_Y\ge0$, vanishing exactly on $Y$, such that near
each point of $Y$, in an adapted chart, $q_Y=x^TB(x,y)x$ with $B$ smooth and
positive definite. Then $\rho_{\mathrm{ff}}=(\beta^*q_Y)^{1/2}$ is a smooth defining
function for the front face; in a polar chart it is $r$ times a positive smooth function.
\item For $H\in\mathcal M_1(X)$ let $H'$ be the closure of $\beta^{-1}(H\setminus Y)$, the
\emph{lift} of $H$. Then
\begin{equation*}
  \beta^*\rho_H=\rho_{\mathrm{ff}}^{\,e_H}\rho_{H'},\qquad e_H=\begin{cases}1,&Y\subset H\\
  0,&Y\not\subset H,\end{cases}
\end{equation*}
with $\rho_{H'}$ a defining function for $H'$. The boundary hypersurfaces of $[X;Y]$
are the components of the front face and of the lifts $H'$; in particular $[X;Y]$ is a
manifold with corners and $\beta$ is an interior b-map. \textup{(}If $Y=H$ is itself a
boundary hypersurface, $\Sphere^{1,0}$ is a point and $[X;H]=X$.\textup{)}
\item Every diffeomorphism $F:X\to X'$ with $F(Y)=Y'$ lifts to a unique diffeomorphism
$[X;Y]\to[X';Y']$, acting on the front face by the projectivized normal differential
$S^+N_pY\ni[v]\mapsto[F_*v]\in S^+N_{F(p)}Y'$.
\item If $U\subset X$ is open, $\beta^{-1}(U)=[U;Y\cap U]$. For manifolds with corners
$X_1$, $X_2$ and $Y_1\subset X_1$ a p-submanifold, $[X_1\times X_2;Y_1\times X_2]=
[X_1;Y_1]\times X_2$.
\end{enumerate}
\end{proposition}
\begin{proof}
(i) Cover $Y$ by adapted charts $U_i$, with normal coordinates $x^{(i)}$, let $U_0=X\setminus
Y$, and let $(\phi_i)$ be a partition of unity subordinate to this cover. Put
$q_Y=\phi_0+\sum_{i\ge1}\phi_i\abs{x^{(i)}}^2$. At a point not in $Y$ some term is
positive. Near $p\in Y$, in an adapted chart with normal coordinates $x$, Taylor's
formula gives $x^{(i)}=A_i(x,y)x$ with $A_i$ invertible (as in the proof of
\cref{blowup:thm:invariance}), so $q_Y=x^TBx$ with
$B=\sum\phi_iA_i^TA_i$ positive definite, $\phi_0$ vanishing near $p$. In polar
coordinates $\beta^*q_Y=r^2\,\omega^TB(r\omega,y)\omega$ with $\omega^TB\omega>0$, so
$(\beta^*q_Y)^{1/2}=r(\omega^TB\omega)^{1/2}$ is smooth; away from the front face
$\beta^*q_Y>0$.

(ii) If $Y\subset H$ then near $p\in Y$ we may take $H=\{x_i=0\}$, $i\le a$, and by
\cref{blowup:lem:bdf} $\rho_H=x_ig$, $g>0$; so $\beta^*\rho_H=r\omega_i\,\beta^*g$, and
$\rho_{H'}=\beta^*\rho_H/\rho_{\mathrm{ff}}$ is a positive multiple of $\omega_i$, the
defining function of the face $\{\omega_i=0\}$ of $\Sphere^{a,b}$, which is $H'$ near
$\mathrm{ff}$. If $Y\not\subset H$ but $p\in H$, then $H$ is $\{y_l=0\}$ for a boundary
coordinate $y_l$, and $\beta^*\rho_H$ is a positive multiple of $y_l$, which is a
coordinate on $[X;Y]$. Away from $\mathrm{ff}$ there is nothing to prove. Thus all the
codimension-one points of $[X;Y]$ lie in $\mathrm{ff}$ or in some $H'$, and each of these
sets has a global defining function. The components of a disconnected front face or
lift are disjoint closed sets, and a defining function for one component is obtained by
replacing $\rho$ by $1$ near the others (using a cut-off).

(iii) Adapted charts are carried by $F$ to adapted charts, and in polar charts the lift
is the map of \cref{blowup:thm:invariance}; by \eqref{blowup:eq:ff-action} it acts on the
front face by the normal differential.

(iv) Both statements follow from the definitions, since adapted charts of $Y_1\subset
X_1$ times charts of $X_2$ are adapted charts of $Y_1\times X_2$.
\end{proof}

\begin{discussion}
The front face remembers \emph{from which direction} a point of $Y$ was approached:
a curve $\gamma$ in $X\setminus Y$ with $\gamma(t)\to p\in Y$ and $\gamma'(0)\notin T_pY$
lifts to a curve ending at the point $[\gamma'(0)]$ of $S^+_pY$. Functions which have
different limits along different directions, such as $x_1/\abs x$, become continuous ---
indeed smooth--- after the blow-up. Functions on $X$ always lift to smooth functions,
so $\beta^*:\CI(X)\to\CI([X;Y])$ is injective, but not surjective.
\end{discussion}

\begin{question}
Why only p-submanifolds? The diagonal of $[0,1]^2$ is a perfectly good submanifold.
\end{question}
\begin{answer}
The construction needs, near each point of $Y$, a chart in which $Y$ is a coordinate
subspace; the polar coordinates in the normal variables are then a manifold with corners,
and \cref{blowup:thm:invariance} makes them independent of the chart. Near $(0,0)$ the
square, in the coordinates $w=x-x'$ (normal to the diagonal) and $x+x'$ (along it), is the
wedge $\{\abs w\le x+x'\}$, not a product of a normal model with a tangential one, and
\cref{blowup:ex:diag-not-p} shows that no chart does better. The standard remedy is to
blow up first the boundary faces where the product structure fails: after blowing up the
corner, the diagonal lifts to a p-submanifold (\cref{blowup:exa:b-double}).
\end{answer}

\subsection{Lifting vector fields}
\begin{proposition}\label{blowup:prop:vf}
Let $Y\subset X$ be a closed p-submanifold. Every smooth vector field on $X$ tangent to
$Y$ lifts to a unique smooth vector field on $[X;Y]$, tangent to the front face. If
$V\in\Vb(X)$ is tangent to $Y$, its lift lies in $\Vb([X;Y])$. The lifts of the elements
of $\Vb(X)$ tangent to $Y$ span $\Vb([X;Y])$ over $\CI([X;Y])$ near the front face.
\end{proposition}
\begin{proof}
Uniqueness follows from density of $[X;Y]\setminus\mathrm{ff}$. Work in an adapted chart,
$x\in\R^{a+b,a}$. A smooth vector field $V=\sum_jc_j\pa_{x_j}+\sum_ld_l\pa_{y_l}$ is tangent
to $Y$ if and only if $c_j(0,y)=0$, i.e.\ $c_j=\sum_ix_ic_{ij}$. The $\pa_{y_l}$ are
unchanged in polar coordinates, and by \eqref{blowup:eq:lin-lift} (which holds verbatim
on $\Sphere^{a,b}$) $x_i\pa_{x_j}$ lifts to $\omega_i\omega_jr\pa_r+\omega_iW_j$, a smooth
vector field tangent to $r=0$. So all such $V$ lift.

If $V\in\Vb(X)$, then by \cref{blowup:lem:Vb} $c_j\in x_j\CI$ for $j\le a$, and $d_l\in
y_l\CI$ for boundary coordinates $y_l$. Hence $V$ is a combination of
\begin{equation}\label{blowup:eq:Vb-Y}
  x_i\pa_{x_k}\ (k>a,\ \text{any } i),\qquad x_m\pa_{x_m}\ (m\le a),\qquad
  \text{and the generators of }\Vb\text{ in }y.
\end{equation}
By \cref{blowup:prop:props}(ii) the boundary hypersurfaces of $[X;Y]$ are the front face
and lifts, and we check tangency in projective coordinates
\eqref{blowup:eq:projective}, which cover the front face. By
\eqref{blowup:eq:proj-lifts}, $x_m\pa_{x_m}$ lifts to $t_m\pa_{t_m}$ ($m\ne j$) or to
$s\pa_s-\sum_rt_r\pa_{t_r}$ ($m=j$), and $x_i\pa_{x_k}$ with $k>a$ lifts to $t_i\pa_{t_k}$,
$\pm\pa_{t_k}$ or $\pm t_i(s\pa_s-\sum_rt_r\pa_{t_r})$; here $t_k$, $k>a$, is an interior
coordinate. All these are in $\Vb$. For spanning, near a point of $\{\sigma\omega_j>0\}$ the
space $\Vb([X;Y])$ is spanned by $s\pa_s$, $t_m\pa_{t_m}$ ($m\le a$, $m\ne j$), $\pa_{t_k}$
($k>a$, $k\ne j$) and the generators in $y$. Now $t_m\pa_{t_m}$ is the lift of
$x_m\pa_{x_m}$; $\pa_{t_k}=\pm$ the lift of $x_j\pa_{x_k}$ (allowed since $k>a$); and
\begin{equation*}
  s\pa_s=\text{lift of }x_j\pa_{x_j}+\sum_{m\le a,m\ne j}t_m\pa_{t_m}+
  \sum_{k>a,k\ne j}t_k\pa_{t_k},
\end{equation*}
where each term on the right is a smooth multiple of a lift.
\end{proof}

\subsection{Vector bundles and tubular neighbourhoods}
For a real vector bundle $E\to M$ of rank $k$ over a manifold with corners, the zero
section $0_E$ is an interior p-submanifold (fibre coordinates are normal coordinates),
and $N_{0_p}0_E=E_p$ canonically. Thus $S^+N0_E=SE=(E\setminus0_E)/\R^+$, the sphere
bundle, and
\begin{equation*}
  [E;0_E]=(E\setminus0_E)\sqcup SE.
\end{equation*}

\begin{proposition}\label{blowup:prop:bundle}
Choose a fibre metric on $E$ and identify $SE$ with the unit sphere bundle. The map
$[0,\infty)\times SE\ni(t,\theta)\mapsto t\theta$ (for $t>0$), $(0,\theta)\mapsto[\theta]$,
is a diffeomorphism onto $[E;0_E]$. Every bundle isomorphism $E\to E'$ covering a
diffeomorphism $M\to M'$ lifts to a diffeomorphism $[E;0_E]\to[E';0_{E'}]$. In
particular, for a real vector space $W$ the blow-up $[W;\{0\}]=(W\setminus0)\sqcup\Sphere W$
is a manifold with boundary diffeomorphic to $[0,\infty)\times\Sphere^{n-1}$, and the action
of $\GL(W)$ on $W$ lifts to a smooth action on it.
\end{proposition}
\begin{proof}
In a local trivialization by a metric-orthonormal frame the map is the identity in polar
coordinates, so it is a diffeomorphism. Bundle maps lift by
\cref{blowup:prop:props}(iii) (in a trivialization they are $F(x,y)=(A(y)x,G(y))$). For
the last part take $M$ a point; the lift of $G\in\GL(W)$ is
$(r,\omega)\mapsto(r\abs{G\omega},G\omega/\abs{G\omega})$, which is smooth jointly in
$(G,r,\omega)$.
\end{proof}

For an interior closed submanifold of a manifold without boundary, the tubular
neighbourhood theorem reduces the general blow-up to this case.

\begin{proposition}\label{blowup:prop:tubular}
Let $X$ be a manifold without boundary, $Y\subset X$ a closed submanifold and $f:D\to D'$
a normal fibration as in \cref{conormal:thm:tubular}. Then $f$ lifts to a diffeomorphism
$[D;Y]\to[D';0_{NY}]\subset[NY;0_{NY}]$ which is the identity on the front face
$SNY$. If $Y$ is compact, a neighbourhood of the front face in $[X;Y]$ is
diffeomorphic to $[0,\epsilon)\times SNY$, with the front face $\{0\}\times SNY$.
\end{proposition}
\begin{proof}
By \cref{blowup:prop:props}(iii,iv) $f$ lifts to a diffeomorphism $[D;Y]\to[D';0_{NY}]$
acting on the front face by the normal differential of $f$, which is the identity by
\cref{conormal:thm:tubular}. By \cref{blowup:prop:bundle}, $[D';0_{NY}]$ is an open
neighbourhood of $\{0\}\times SNY$ in $[0,\infty)\times SNY$, which contains
$[0,\epsilon)\times SNY$ if $Y$ is compact.
\end{proof}

%======================================================================
\section{Functions and densities on blown-up spaces}\label{blowup:sec:functions}

\begin{lemma}\label{blowup:lem:flat}
Let $Q\subset\R^{d,c}$ be open and $G\in\CI((0,\tau_0)\times Q)$. Suppose that for every
compact $K\subset Q$, every derivative $\pa^\alpha G$ and every $N$ there is $C$ with
$\abs{\pa^\alpha G(\tau,q)}\le C\tau^N$ for $q\in K$. Then $G$, extended by $0$ to
$\tau=0$, lies in $\CI([0,\tau_0)\times Q)$ and vanishes to infinite order at
$\tau=0$.
\end{lemma}
\begin{proof}
The extension by $0$ to $\tau\le0$ has all derivatives bounded on compact subsets of the
interior of $(-\infty,\tau_0)\times Q$, and each derivative tends to $0$ as $\tau\downarrow0$;
by the mean value theorem the derivatives of the extension are continuous across
$\tau=0$ and vanish there. So it is smooth by \cref{blowup:rem:smooth}.
\end{proof}

\begin{example}\label{blowup:exa:homogeneous}
If $f\in\CI(\R^n\setminus0)$ is homogeneous of degree $\mu\in\C$, then on $[\R^n;\{0\}]$,
$\beta^*f=r^\mu f(\omega)$ is $r^\mu$ times a smooth function. Thus $\abs z=r$ and
$z_1/\abs z=\omega_1$ are smooth on $[\R^n;\{0\}]$, though not on $\R^n$, and the terms of
the expansions in \cref{conormal:prop:homogeneous} are powers of $r$ times smooth
functions.
\end{example}

\begin{proposition}\label{blowup:prop:density}
Let $Y\subset X$ be a closed p-submanifold of codimension $k$ and $\nu$ a smooth positive
density on $X$. Then $\beta^*\nu=\rho_{\mathrm{ff}}^{\,k-1}\tilde\nu$ with $\tilde\nu$ a smooth
positive density on $[X;Y]$. In polar coordinates $\beta^*\abs{dx\,dy}=r^{k-1}\abs{dr\,d\omega\,dy}$,
where $d\omega$ is the Euclidean surface measure on $\Sphere^{a,b}$, and in projective
coordinates $\beta^*\abs{dx\,dy}=s^{k-1}\abs{ds\,dt\,dy}$.
\end{proposition}
\begin{proof}
The polar formula is the classical change of variables $x=r\omega$. In projective
coordinates $x_j=\sigma s$, $x_k=st_k$, the Jacobian matrix of $(s,t)\mapsto x$ is
triangular after moving the $j$th row to the top, with determinant $\pm s^{k-1}$. Since
$r$ and $s$ are positive multiples of $\rho_{\mathrm{ff}}$, the claim follows.
\end{proof}

\subsection{The semiclassical double space}
The following result describes exactly when a family of kernels which concentrates at
the diagonal as $\eps\downarrow0$ becomes smooth after blowing up $\{0\}\times\Diag$. It is
the geometric content of \cref{semiclassical:rem:double-space} and
\cref{semiclassical:prop:smoothing-X}(iii). We allow a manifold $P$ (without boundary)
of parameters, which gives the adiabatic case.

Let $\Omega\subset\R^n$ be open and set
\begin{equation*}
  \mathcal X=[0,1]_\eps\times\Omega_x\times\Omega_y\times P,\quad
  \mathcal Y=\{0\}\times\Diag_\Omega\times P,\quad
  \mathcal D=[\mathcal X;\mathcal Y].
\end{equation*}
In the coordinates $(\eps,x,w=x-y,p)$, $\mathcal Y=\{\eps=0,\ w=0\}$ is a p-submanifold of
normal type $(1,n)$, so the front face is a bundle of closed half-spheres $\Sphere^{1,n}$
over $\Diag_\Omega\times P$. Besides the front face and the lift of $\{\eps=1\}$, $\mathcal D$
has one more boundary hypersurface, the lift $\mathrm{sf}$ of $\{\eps=0\}$ (the \emph{side
face}).

\begin{question}
Why blow up, instead of simply using the rescaled variable $Z=(x-y)/\eps$?
\end{question}
\begin{answer}
Because a kernel lives in two regimes at once. Near the diagonal, with $Z$ bounded, it
is $K(\eps,x,Z)$; away from the diagonal, with $x\ne y$ fixed and $\eps\to0$, it should
vanish rapidly. The coordinate $Z$ covers only the first regime, and the coordinates
$(\eps,x,y)$ only the second. The blown-up space contains both as open sets, glued
along the region where $\abs Z$ is large; rapid decay in $Z$ becomes vanishing to infinite
order at the side face. Moreover the blown-up space is defined without coordinates, which
makes the description invariant (\cref{blowup:cor:sc-manifold}).
\end{answer}

\begin{proposition}\label{blowup:prop:sc-kernels}
\begin{enumerate}[label=\textup{(\roman*)}]
\item The map $(\eps,x,Z,p)\mapsto\beta^{-1}(\eps,x,x-\eps Z,p)$, $\eps>0$, extends to a
diffeomorphism $\Psi$ from $\{(\eps,x,Z,p)\in[0,1]\times\Omega\times\R^n\times P;\
x-\eps Z\in\Omega\}$ onto $\mathcal D\setminus\mathrm{sf}$, mapping $\{\eps=0\}$ onto
$\mathrm{ff}\setminus\mathrm{sf}$. Thus the interior of the front face is
$\Omega\times\R^n\times P$, with $Z$ a tangent vector at $x$.
\item Let $K\in\CI([0,1]\times\Omega\times\R^n\times P)$ satisfy, for every compact
$L\subset\Omega\times P$, all derivatives (in all variables) and all $N$,
\begin{equation}\label{blowup:eq:K-est}
  \sup_{\eps\in[0,1],\ (x,p)\in L,\ Z\in\R^n}\ang Z^N\abs{\pa^\alpha K(\eps,x,Z,p)}<\infty.
\end{equation}
Then $k(\eps,x,y,p)=K(\eps,x,(x-y)/\eps,p)$, $\eps>0$, is such that $\beta^*k$ extends to an
element of $\CI(\mathcal D)$ vanishing to infinite order at $\mathrm{sf}$.
\item Conversely, let $k\in\CI((0,1]\times\Omega^2\times P)$ be supported in
$(0,1]\times L$ for a compact $L\subset\Omega^2\times P$, and suppose $\beta^*k$ extends to an
element of $\CI(\mathcal D)$ vanishing to infinite order at $\mathrm{sf}$. Then
$K(\eps,x,Z,p)=k(\eps,x,x-\eps Z,p)$ \textup{(}$:=0$ if $x-\eps Z\notin\Omega$\textup{)}
extends to an element of $\CI([0,1]\times\Omega\times\R^n\times P)$ satisfying
\eqref{blowup:eq:K-est}.
\end{enumerate}
\end{proposition}
\begin{proof}
The roadmap: (i) is the projective chart with $\eps$ as the projective variable; it
covers everything except the side face. Near the side face we use projective charts
based on a component $w_j$ of $w$, in which large $\abs Z$ corresponds to the side face
defining function going to $0$; rapid decay in $Z$ then translates into vanishing to
infinite order at $\mathrm{sf}$, and conversely.

(i) In polar coordinates around $\mathcal Y$ the normal variables are $(\eps,w)\in
[0,\infty)\times\R^n$, with $\omega=(\omega_0,\omega')\in\Sphere^{1,n}$. The side face is
the closure of $\beta^{-1}(\{\eps=0,\ w\ne0\})$, which is $\{\omega_0=0\}$. On
$\{\omega_0>0\}=\mathcal D\setminus\mathrm{sf}$ the projective coordinates of
\cref{blowup:lem:projective} with $j=0$ are $s=\eps$, $t=w/\eps=Z$, together with $x$,
$p$. This is $\Psi^{-1}$.

(ii) By (i), $\beta^*k=K\circ\Psi^{-1}$ is smooth on $\mathcal D\setminus\mathrm{sf}$. A point
of $\mathrm{sf}$ has $\omega_0=0$, hence $\omega_j\ne0$ for some $j\ge1$; say $\omega_j>0$
(the case $\omega_j<0$ is the same with signs changed). Near it we use the projective
coordinates
\begin{equation*}
  s=w_j\ge0,\quad\tau=\eps/w_j\ge0,\quad t_k=w_k/w_j\ (k\ne0,j),\quad x,\ p,
\end{equation*}
in which $\mathrm{ff}=\{s=0\}$ and $\mathrm{sf}=\{\tau=0\}$. With $\theta\in\R^n$ having $1$ in slot $j$
and $t_k$ in slot $k$, we have $\eps=\tau s$ and $Z=w/\eps=\theta/\tau$, so
\begin{equation*}
  \beta^*k=G(\tau,s,t,x,p)=K\big(\tau s,x,\theta/\tau,p\big),\qquad\tau>0.
\end{equation*}
Since $\pa_\tau$, $\pa_s$, $\pa_{t_k}$ of the arguments are $(s,-\theta/\tau^2)$,
$(\tau,0)$ and $(0,e_k/\tau)$, induction shows that every derivative of $G$ is a finite
sum of terms $c(s,t)\tau^{-q}(\pa^\gamma K)(\tau s,x,\theta/\tau,p)$ with $c$ polynomial.
On a compact set of $(s,t,x,p)$, $\abs{\theta/\tau}\ge1/\tau$, so by \eqref{blowup:eq:K-est}
each term is $O(\tau^{N-q})$ for every $N$. By \cref{blowup:lem:flat}, $G$ extends
smoothly to $\tau=0$ and vanishes to infinite order there.

(iii) By (i), $K=\beta^*k\circ\Psi$ is smooth; it vanishes near points with
$x-\eps Z\in\pa\Omega$ because $L$ is compact, so its extension by $0$ is smooth. For the
estimates, consider $\abs Z\ge2\sqrt n$ and choose $j$ with $\abs{Z_j}=\max_k\abs{Z_k}$, say
$Z_j>0$. In the chart of (ii), $\tau=1/Z_j$, $s=\eps Z_j$, $t=Z/Z_j$ with $\abs{t_k}\le1$
and $\tau\le1$. The set $S=\beta^{-1}([0,1]\times L)$ is compact, and its intersection with
$\{\tau\le2,\ \abs{t}_\infty\le2\}$ is a compact subset $S'$ of the chart (there all
components of $\omega$ are at most $2\omega_j$, so $\omega_j\ge(4n+4)^{-1/2}$). Decreasing
$\tau$ with $(s,t,x,p)$ fixed decreases $\eps=\tau s$ and fixes $(x,y,p)$, so $S'$ contains
with each point the segment joining it to $\{\tau=0\}$. Since $G$ vanishes to infinite order
at $\tau=0$, Taylor's formula in $\tau$ gives $\abs{\pa^\gamma G}\le C_{\gamma,N}\tau^N$ on $S'$,
and $G=0$ off $S$. Now
$K(\eps,x,Z,p)=G(1/Z_j,\eps Z_j,Z/Z_j,x,p)$, and each derivative of order $q$ of the maps
$(\eps,Z)\mapsto(1/Z_j,\eps Z_j,Z/Z_j)$ is bounded by $C\ang Z^q$ on this region. By the
chain rule $\abs{\pa^\alpha K}\le C\ang Z^{\abs\alpha}\sup_{\abs\gamma\le\abs\alpha}
\abs{\pa^\gamma G}\le C'\ang Z^{\abs\alpha}Z_j^{-N}\le C''\ang Z^{\abs\alpha-N}$, since
$Z_j\ge\abs Z/\sqrt n$. As $N$ is arbitrary, \eqref{blowup:eq:K-est} holds.
\end{proof}

\begin{figure}[ht]
\centering
\begin{tikzpicture}[>=Stealth,scale=1.0]
  \begin{scope}
    \fill[gray!12] (0,0) rectangle (2.6,2.6);
    \draw[very thick] (0,2.6)--(0,0)--(2.6,0);
    \draw[thick,dashed] (0,0)--(2.6,2.6);
    \fill (0,0) circle (2pt);
    \node[font=\small] at (2.1,1.5) {$\Diag$};
    \node[below,font=\small] at (1.3,0) {$x'=0$};
    \node[left,font=\small] at (0,1.3) {$x=0$};
    \node[font=\small] at (1.3,-0.9) {$[0,\infty)^2$};
  \end{scope}
  \draw[<-] (3.1,1.3)--(4.3,1.3) node[midway,above,font=\small]{$\beta$};
  \begin{scope}[xshift=4.9cm]
    \fill[gray!12] (0,0) rectangle (2.6,2.6);
    \fill[white] (0,0) -- (1,0) arc (0:90:1) -- cycle;
    \draw[very thick] (0,2.6)--(0,1) arc (90:0:1) -- (2.6,0);
    \draw[thick,dashed] (45:1)--(2.6,2.6);
    \fill (45:1) circle (1.6pt);
    \node[font=\small] at (0.35,0.3) {ff};
    \node[font=\small] at (2.1,1.5) {$\Diag_{\mathrm b}$};
    \foreach \a in {14,27,63,76} \draw[gray] (\a:1)--(\a:2.45);
    \node[font=\footnotesize,align=left] at (3.55,1.0) {gray: $t=x'/x$\\ constant};
    \node[font=\small] at (1.3,-0.9) {$[[0,\infty)^2;\{0\}]$};
  \end{scope}
\end{tikzpicture}
\caption{Blowing up the corner of $[0,\infty)^2$. The front face is a quarter circle
$\Sphere^{2,0}$. The diagonal, which is not a p-submanifold, lifts to the interior
p-submanifold $\Diag_{\mathrm b}=\{t=1\}$ meeting the front face transversally. In the
projective coordinates $s=x$, $t=x'/x$ (valid away from the left edge), the front face
is $\{s=0\}$ and the lifted edge $\{x'=0\}$ is $\{t=0\}$.}
\label{blowup:fig:corner}
\end{figure}

\begin{corollary}\label{blowup:cor:sc-manifold}
Let $X$ be a compact manifold of dimension $n$, $\mathcal D_X=[[0,1]\times X^2;\{0\}\times\Diag]$,
and $k\in\CI((0,1]\times X^2)$. Then $\beta^*k$ extends to an element of $\CI(\mathcal
D_X)$ vanishing to infinite order at the lift of $\{\eps=0\}$ if and only if
\textup{(a)} $k$ extends smoothly to $[0,1]\times(X^2\setminus\Diag)$, vanishing to
infinite order at $\eps=0$, and \textup{(b)} for every coordinate chart $W\cong\Omega
\subset\R^n$ and $\chi\in\CIc(W\times W)$, $\chi k=K_\chi(\eps,x,(x-y)/\eps)$ with $K_\chi$
as in \cref{blowup:prop:sc-kernels}(ii) \textup{(}$P$ a point\textup{)}. The front face of
$\mathcal D_X$ is the bundle over $\Diag\cong X$ of closed half-spheres
$S^+N_{(0,x,x)}(\{0\}\times\Diag)\subset S(\R\times T_xX)$, and its interior is $TX$ by
$[\eps:v]\mapsto v/\eps$.
\end{corollary}
\begin{proof}
On $[0,1]\times(X^2\setminus\Diag)$, $\beta$ is a diffeomorphism and the lift of
$\{\eps=0\}$ is $\{\eps=0\}$, so the first condition there is (a). Near $\Diag$ use
\cref{blowup:prop:props}(iv): $\beta^{-1}([0,1]\times W^2)=[[0,1]\times W^2;\{0\}\times\Diag_W]$.
If $\beta^*k$ has the property, so does $\beta^*(\chi k)=(\beta^*\chi)(\beta^*k)$, and
\cref{blowup:prop:sc-kernels}(iii) gives (b). Conversely, write $k=\chi_0k+\sum_i\chi_ik$
with $\chi_i\in\CIc(W_i^2)$ and $\supp\chi_0\cap\Diag=\emptyset$; the first term is handled by
(a), the others by (b) and \cref{blowup:prop:sc-kernels}(ii). The normal bundle of
$\{0\}\times\Diag$ is $\R\times N\Diag$, with $N_{(x,x)}\Diag\cong T_xX$ by
$(v,v')\mapsto v-v'$, and the inward part is $\{\eps\ge0\}$; in the coordinates of
\cref{blowup:prop:sc-kernels}(i) the point $[\eps:v]$ with $\eps>0$ has $Z=v/\eps$.
\end{proof}

\begin{remark}\label{blowup:rem:adiabatic}
The same argument applies to the adiabatic families of \cref{semiclassical:sec:adiabatic}.
For a fibration $\phi:Y\to Z$ with compact fibres, the fibre product
$Y^2_\phi=\{(y,y');\ \phi(y)=\phi(y')\}=(\phi\times\phi)^{-1}(\Diag_Z)$ is a closed submanifold of
$Y^2$ (since $\phi\times\phi$ is a submersion), and the adiabatic double space is
$[[0,1]\times Y^2;\{0\}\times Y^2_\phi]$. In fibred coordinates $(z,\tilde z)$, with
$\tilde z$ the base coordinates, $Y^2_\phi=\{\tilde z=\tilde z'\}$, so this is the space
$\mathcal D$ above with $\Omega$ a base chart and parameters $p=(z,z')$. Thus, by
\cref{blowup:prop:sc-kernels}, kernels $B(\delta,z,z',\tilde z,(\tilde z-\tilde z')/\delta)$
with $B\in\mathcal K_{\mathrm{ad}}$ (\cref{semiclassical:def:adiabatic}), cut off to compact
sets, are exactly the smooth functions on the adiabatic double space vanishing to infinite
order at the side face. The interior of its front face is the bundle
$Y^2_\phi\times_ZTZ$, on which the adiabatic symbol is the fibrewise Fourier transform.
\end{remark}

%======================================================================
\section{Lifts of submanifolds and commutation of blow-ups}\label{blowup:sec:commute}

To blow up twice we must know what the second centre becomes after the first blow-up.

\begin{definition}\label{blowup:def:lift}
Let $\beta:[X;Y]\to X$ and $Z\subset X$ closed. If $Z\subset Y$, the \emph{lift} of $Z$ is
$\beta^{-1}(Z)$. If $Z\setminus Y$ is dense in $Z$, the lift is
$\beta^\#Z=\overline{\beta^{-1}(Z\setminus Y)}$. We write $[X;Y;Z]=[[X;Y];\text{lift of }Z]$
when the lift is a p-submanifold.
\end{definition}

\begin{proposition}\label{blowup:prop:lift}
Let $Y,Z\subset X$ be closed joint p-submanifolds.
\begin{enumerate}[label=\textup{(\roman*)}]
\item If $Z\subset Y$, then $\beta^{-1}(Z)=S^+NY|_Z$ is a p-submanifold of $[X;Y]$
contained in the front face.
\item If $Z\setminus Y$ is dense in $Z$, then $\beta^\#Z$ is a p-submanifold of $[X;Y]$,
and $\beta$ restricts to the blow-down map of $[Z;Y\cap Z]$, i.e.\ $\beta^\#Z=[Z;Y\cap Z]$.
In particular, if $Y\subset Z$ then $\beta^\#Z=[Z;Y]$, and if $Y\cap Z=\emptyset$ then
$\beta^\#Z=\beta^{-1}(Z)\cong Z$.
\end{enumerate}
\end{proposition}
\begin{proof}
Work in a joint chart $(z_1,\dots,z_n)$ with $Y=\{z_A=0\}$, $Z=\{z_B=0\}$. Then $z_A$ are
normal coordinates for $Y$, and the other $z_i$ are tangential.

(i) Here $A\subset B$, and in polar coordinates $\beta^{-1}(Z)=\{r=0,\ z_{B\setminus A}=0\}$,
a p-submanifold.

(ii) Density of $Z\setminus Y$ means $A\not\subset B$. For $r>0$,
$\beta^{-1}(Z\setminus Y)=\{\omega_{A\cap B}=0,\ z_{B\setminus A}=0\}$. A point of the closure
on the front face has $\omega_{A\cap B}=0$, so $\omega_j\ne0$ for some $j\in A\setminus B$,
and conversely a point of $\{r=0,\ \omega_{A\cap B}=0,\ z_{B\setminus A}=0\}$ is a limit of points
with $r>0$. In projective coordinates based at such a $j$, with $s=\abs{z_j}$ and
$t_k=z_k/\abs{z_j}$, the closure is
\begin{equation*}
  \beta^\#Z=\{t_{A\cap B}=0,\ z_{B\setminus A}=0\},
\end{equation*}
a p-submanifold. The remaining coordinates $(s,t_{A\setminus(B\cup\{j\})},z_{\text{rest}})$ on
$\beta^\#Z$ are exactly projective coordinates for the blow-up of $Y\cap Z=\{z_{A\cup B}=0\}$
in $Z$, whose normal coordinates in $Z$ are $z_{A\setminus B}$; so $\beta|_{\beta^\#Z}$ is the
blow-down map of $[Z;Y\cap Z]$ in each chart.
\end{proof}

\begin{proposition}\label{blowup:prop:commute}
Let $Y,Z\subset X$ be closed joint p-submanifolds such that one of the following holds:
$Y\cap Z=\emptyset$; $Y$ and $Z$ are transversal; $Y\subset Z$; or $Z\subset Y$. Then the
lifts in \cref{blowup:def:lift} are p-submanifolds, and the identity of $X\setminus(Y\cup Z)$
extends to a diffeomorphism $[X;Y;Z]\cong[X;Z;Y]$ commuting with the blow-down maps to
$X$.
\end{proposition}
\begin{proof}
The lifts are p-submanifolds by \cref{blowup:prop:lift}. If a diffeomorphism extending
the identity exists over each set of an open cover of $X$, the local extensions agree on
overlaps (the interior is dense) and define a smooth map $[X;Y;Z]\to[X;Z;Y]$; the same
holds for the inverse, and the composites are the identity on a dense set. So the
statement is local on $X$, and it suffices to work in a joint chart. Away from $Y\cap Z$
only one blow-up has an effect. Near a point of $Y\cap Z$ we treat the cases in turn.

\emph{Transversal.} Coordinates $(u,v,w)$ with $Y=\{u=0\}$, $Z=\{v=0\}$. By
\cref{blowup:prop:props}(iv) both iterated blow-ups are
$[0,\infty)\times\Sphere_u\times[0,\infty)\times\Sphere_v\times W$, with the interior
map $(u,v,w)\mapsto(\abs u,u/\abs u,\abs v,v/\abs v,w)$; here $\Sphere_u$, $\Sphere_v$ are
the corner spheres of the normal models.

\emph{Nested.} Let $Y\subset Z$, with $Z=\{u=0\}$, $Y=\{u=0,\ v=0\}$, $u\in\R^{a_1+b_1,a_1}$,
$v\in\R^{a_2+b_2,a_2}$, and let $H=\Sphere^{1+a_2,b_2}=\{(\alpha,\phi);\ \alpha\ge0,\ \phi\in
\R^{a_2+b_2,a_2},\ \alpha^2+\abs\phi^2=1\}$. We show that both spaces are identified with
\begin{equation*}
  [0,\infty)\times H\times\Sphere^{a_1,b_1}\times W\quad\text{by}\quad
  (u,v,w)\longmapsto\Big(\abs{(u,v)},\ \frac{(\abs u,v)}{\abs{(u,v)}},\ \frac u{\abs u},\ w\Big).
\end{equation*}
First blow up $Z$: $[X;Z]=[0,\infty)_\rho\times\Sphere^{a_1,b_1}_\theta\times\R_v\times W$ with
$\rho=\abs u$, and $\beta_Z^{-1}(Y)=\{\rho=0,\ v=0\}$ has normal variables $(\rho,v)\in
\R^{1+a_2+b_2,1+a_2}$. Blowing it up introduces $\sigma=\abs{(\rho,v)}=\abs{(u,v)}$ and
$(\rho,v)/\sigma\in H$, as claimed. Now blow up $Y$ first: $[X;Y]=[0,\infty)_R\times
\Sphere^{a_1+a_2,b_1+b_2}\times W$ with $R=\abs{(u,v)}$ and $(\mu,\nu)=(u,v)/R$, and
$\beta^\#Z=\{\mu=0\}$. It suffices to show that $(\mu,\nu)\mapsto((\abs\mu,\nu),\mu/\abs\mu)$
extends to a diffeomorphism $[\Sphere^{a_1+a_2,b_1+b_2};\{\mu=0\}]\to H\times
\Sphere^{a_1,b_1}$. Away from $\mu=0$ it is a diffeomorphism onto $\{\alpha>0\}$, with
inverse $((\alpha,\phi),\theta)\mapsto(\alpha\theta,\phi)$. Near a point with $\mu=0$ some
$\nu_m\ne0$, and $(\mu,\nu')$, $\nu'=(\nu_k)_{k\ne m}$, are coordinates on the sphere with
$\{\mu=0\}$ a coordinate p-submanifold; polar coordinates around it are
$(\abs\mu,\mu/\abs\mu,\nu')$. Near the image point, $(\alpha,\phi')$ are coordinates on $H$,
and in these coordinates the map is $(\abs\mu,\mu/\abs\mu,\nu')\mapsto((\abs\mu,\nu'),
\mu/\abs\mu)$, the identity. The case $Z\subset Y$ is the same with the names exchanged.
\end{proof}

\begin{discussion}
The proof shows what goes wrong in general: for joint $Y$ and $Z$ the lifts are
p-submanifolds whatever their relative position, but the two iterated blow-ups introduce different
``rates'' at which points approach $Y\cap Z$. In \cref{blowup:ex:noncommute} the function
$z_2/\abs{(z_1,z_2)}$ is continuous on one of the two spaces but not on the other. The
commonest cure is to blow up $Y\cap Z$ first.
\end{discussion}

\begin{corollary}\label{blowup:cor:intersection}
Let $Y,Z$ be closed joint p-submanifolds, and $W=Y\cap Z$, with $Y\setminus W$ dense in
$Y$ and $Z\setminus W$ dense in $Z$ \textup{(}$W$ is a p-submanifold, by the joint
charts\textup{)}. Then the lifts of $Y$ and $Z$ to $[X;W]$ are disjoint
p-submanifolds, so $[X;W;Y;Z]\cong[X;W;Z;Y]$.
\end{corollary}
\begin{proof}
In a joint chart, $W=\{z_{A\cup B}=0\}$, and in polar coordinates around $W$ the lifts are
$\{\omega_A=0\}$ and $\{\omega_B=0\}$ (\cref{blowup:prop:lift}(ii)), which cannot both hold on
the sphere in $z_{A\cup B}$. Off the front face they are $Y\setminus W$ and $Z\setminus W$,
which are disjoint.
\end{proof}

\subsection{Relative compactifications}
We now prove the identifications used in \cref{ch:distributions,ch:product}. Let
$W=U\oplus V$ be real vector spaces with Euclidean norms, of dimensions $k,l\ge1$, with
elements $w=(\zeta,\eta)$. The relative compactification of $W$ with respect to $V$
(\eqref{distributions:eq:relRC}, and \cref{ch:product}) is
\begin{equation*}
  \relRC VW=\RC U\times\RC V,\qquad W\ni(\zeta,\eta)\longmapsto(\zeta,q)\in U\times V,\quad
  q=\eta/\ang\zeta,
\end{equation*}
with boundary hypersurfaces $H_W=\pa\RC U\times\RC V$ and $H_V=\RC U\times\pa\RC V$, and
$\Sphere V=\pa\RC V\subset\pa\RC W=\Sphere W$ is the sphere at infinity of $V$.

\begin{theorem}\label{blowup:thm:relative}
$\Sphere V$ is a p-submanifold of $\RC W$, of codimension $k+1$, and the identity of $W$
extends to a diffeomorphism $\relRC VW\to[\RC W;\Sphere V]$ whose composite with the
blow-down map is the map $\varpi$ of \cref{product:prop:blow-down}. It maps $H_V$ onto
the front face and $H_W$ onto the lift of $\Sphere W$, which is $[\Sphere W;\Sphere V]$.
\end{theorem}
\begin{proof}
\emph{Coordinates near $\Sphere V$.} By \cref{distributions:lem:collar}, $(1/\abs w,w/\abs
w)$ are coordinates on $\RC W$ near its boundary. Where the $V$-component of $w/\abs w$
is non-zero, we may replace them by
\begin{equation*}
  x=1/\abs\eta,\qquad v=\zeta/\abs\eta\in U,\qquad\theta=\eta/\abs\eta\in\Sphere V,
\end{equation*}
since $x=s/\abs{\omega_V}$, $v=\omega_U/\abs{\omega_V}$, $\theta=\omega_V/\abs{\omega_V}$ in
terms of $(s,\omega)=(1/\abs w,w/\abs w)$, with smooth inverse
$\abs{\omega_V}=(1+\abs v^2)^{-1/2}$, $\omega_U=v\abs{\omega_V}$, $\omega_V=\theta\abs{\omega_V}$,
$s=x\abs{\omega_V}$. In these coordinates $\Sphere V=\{x=0,\ v=0\}$, a p-submanifold with
normal coordinates $(x,v)\in\R^{1+k,1}$.

\emph{Near the front face.} Polar coordinates around $\Sphere V$ are
\begin{equation*}
  R=\abs{(x,v)}=\frac{\ang\zeta}{\abs\eta}=\frac1{\abs q},\qquad
  \frac{(x,v)}R=\Big(\frac1{\ang\zeta},\frac\zeta{\ang\zeta}\Big)=\mathrm{RC}_U(\zeta),\qquad
  \theta=\frac{\eta}{\abs\eta}=\frac q{\abs q}.
\end{equation*}
The second is a point of the closed half-sphere $\Sphere^{1,k}$, which is $\RC U$
(\cref{distributions:def:RC}). On the other hand, $(1/\abs q,q/\abs q)$ are collar
coordinates on $\RC V$ near $\pa\RC V$. Hence, on the sets $\{R<\delta\}$ (which lie in the
coordinate domains for $\delta<1$, since $R<\delta$ forces $\abs\eta>\delta^{-1}\ang\zeta>1$), the identity
of $W$ is the identity map $(R,h,\theta)\mapsto(R,h,\theta)$ of $[0,\delta)\times\Sphere^{1,k}
\times\Sphere V$, which is a diffeomorphism of a neighbourhood of $H_V$ onto a neighbourhood
of the front face, mapping $H_V=\{1/\abs q=0\}$ to $\{R=0\}$.

\emph{Away from the front face.} On $\RC U\times\RC V\setminus H_V$ we show that $\varpi$
is a diffeomorphism onto $\RC W\setminus\Sphere V=[\RC W;\Sphere V]\setminus\mathrm{ff}$. In the
interior it is the identity. Near a point of $H_W\setminus H_V$ use the coordinates
$(x',\theta',v')=(1/\abs\zeta,\zeta/\abs\zeta,\eta/\abs\zeta)$ on $\RC W$ near $\Sphere W\setminus\Sphere V$
(the analogue of $(x,v,\theta)$ with $U$ and $V$ exchanged) and $(s_U,\theta',q)$ on
$\RC U\times V$, with $s_U=1/\abs\zeta$. Since $q=\eta/\ang\zeta=v'(1+x'^2)^{-1/2}$, the
map is $(x',\theta',v')\mapsto(x',\theta',v'(1+x'^2)^{-1/2})$, a diffeomorphism. It is bijective
from $H_W\setminus H_V$ onto $\Sphere W\setminus\Sphere V$ (both are parametrized by
$(\theta',v')\in\Sphere U\times V$), hence a diffeomorphism.

The two maps agree on the interior, so they define a diffeomorphism of $\relRC VW$ onto
$[\RC W;\Sphere V]$. It maps $H_W$, the closure of $H_W\setminus H_V$, onto the closure of the
lift of $\Sphere W\setminus\Sphere V$, which is $[\Sphere W;\Sphere V]$ by
\cref{blowup:prop:lift}(ii).
\end{proof}

The main face $H_W=[\Sphere W;\Sphere V]$ fibres over $\Sphere(W/V)$; this is the map
$\gamma$ used in \cref{ch:families}.

\begin{lemma}\label{blowup:lem:sphere-fibration}
Let $V\subset W$ be a subspace, $\pi:W\to W/V$ the quotient map and $U$ a complement of $V$
with Euclidean norms on $U$, $V$. The map $\Sphere W\setminus\Sphere V\ni[w]\mapsto[\pi w]\in\Sphere(W/V)$
extends to a smooth map $\gamma:[\Sphere W;\Sphere V]\to\Sphere(W/V)$, which is a trivial
fibre bundle with fibre $\RC V$: the map $[(\zeta,\eta)]\mapsto(\zeta/\abs\zeta,\eta/\abs\zeta)$
extends to a diffeomorphism $[\Sphere W;\Sphere V]\cong\Sphere U\times\RC V$. The extension
$\gamma$ depends only on $V\subset W$. The same holds fibrewise for a vector bundle $W$
with a subbundle $V$, giving a fibre bundle $[\Sphere W;\Sphere V]\to\Sphere(W/V)$ with
fibres $\RC V$.
\end{lemma}
\begin{proof}
Realize $\Sphere W$ as the unit sphere $\abs\zeta^2+\abs\eta^2=1$. On the complement of
$\Sphere V=\{\zeta=0\}$ the map $(\zeta,\eta)\mapsto(\zeta/\abs\zeta,\eta/\abs\zeta)\in\Sphere U\times V$
is a diffeomorphism, with inverse $(\vartheta,p)\mapsto(\vartheta,p)/\abs{(\vartheta,p)}$. Near
$\Sphere V$ the functions $(\zeta,\theta=\eta/\abs\eta)$ are coordinates on $\Sphere W$ (with
$\abs\eta=(1-\abs\zeta^2)^{1/2}$), and $\Sphere V=\{\zeta=0\}$; polar coordinates are
$(a,\vartheta,\theta)=(\abs\zeta,\zeta/\abs\zeta,\theta)$. Near $\Sphere U\times\pa\RC V$ we use the
collar coordinates $(\vartheta,1/\abs p,p/\abs p)$ with $p=\eta/\abs\zeta$, and in these
coordinates the map is $(a,\vartheta,\theta)\mapsto(\vartheta,a(1-a^2)^{-1/2},\theta)$, a
diffeomorphism for $0\le a<1$. Composing with the projection to $\Sphere U\cong\Sphere(W/V)$
gives $\gamma$, since $\pi(\zeta,\eta)=\zeta$. The map $[w]\mapsto[\pi w]$ is intrinsic on the
dense set $\Sphere W\setminus\Sphere V$, so its continuous extension is unique. For bundles, use
local trivializations of $W$ in which $V$ is a coordinate subbundle and fibre metrics; the
construction depends smoothly on parameters.
\end{proof}

\begin{figure}[ht]
\centering
\begin{tikzpicture}[>=Stealth,scale=0.9]
  \begin{scope}
    \fill[gray!12] (0,0) circle (1.5);
    \draw[very thick] (0,0) circle (1.5);
    \foreach \a in {0,90,180,270} \fill (\a:1.5) circle (2pt);
    \node[right,font=\footnotesize] at (0:1.5) {$\Sphere U$};
    \node[left,font=\footnotesize] at (180:1.5) {$\Sphere U$};
    \node[above,font=\footnotesize] at (90:1.5) {$\Sphere V$};
    \node[below,font=\footnotesize] at (270:1.5) {$\Sphere V$};
    \node[font=\small] at (0,-2.3) {$\RC{W}$};
  \end{scope}
  \draw[<-] (2.3,0)--(3.3,0);
  \begin{scope}[xshift=5.4cm]
    \fill[gray!12] (22.5:1.6)--(67.5:1.6)--(112.5:1.6)--(157.5:1.6)--(202.5:1.6)
      --(247.5:1.6)--(292.5:1.6)--(337.5:1.6)--cycle;
    \foreach \c in {0,90,180,270}
      \draw[very thick,densely dashed] (\c-22.5:1.6)--(\c+22.5:1.6);
    \foreach \c in {45,135,225,315}
      \draw[very thick] (\c-22.5:1.6)--(\c+22.5:1.6);
    \node[font=\footnotesize] at (0,0) {$\relRC{U,V}{W}$};
    \node[font=\small] at (0,-2.3) {octagon};
  \end{scope}
  \draw[->] (7.5,0)--(8.5,0);
  \begin{scope}[xshift=10.4cm]
    \fill[gray!12] (-1.2,-1.2) rectangle (1.2,1.2);
    \draw[very thick] (-1.2,-1.2) rectangle (1.2,1.2);
    \foreach \x in {-1.2,1.2} \foreach \y in {-1.2,1.2} \fill (\x,\y) circle (2pt);
    \node[font=\small] at (0,-2.3) {$\RC V\times\RC U$};
  \end{scope}
\end{tikzpicture}
\caption{The doubly-relative compactification for $\dim U=\dim V=1$. Blowing up the four
points $\pm\Sphere U$, $\pm\Sphere V$ of the circle at infinity of $\RC W$ (left), or the four
corners $\Sphere V\times\Sphere U$ of the square (right), gives the same octagon. Its dashed
sides are the front faces of the left-hand blow-up and the remains of the sides of the
square; its solid sides are the remains of the circle and the front faces of the
right-hand blow-up.}
\label{blowup:fig:octagon}
\end{figure}

\begin{proposition}\label{blowup:prop:doubly}
Let $W=V\times U$ with $\dim U,\dim V\ge1$, and let
$E:W\to\R^4\times V\times W\times U$ be the embedding of \cref{distributions:ex:doubly}.
Then $E$ extends to smooth embeddings of both
\begin{equation*}
  \mathcal D_1=[\RC W;\Sphere U,\Sphere V]\quad\text{and}\quad
  \mathcal D_2=[\RC V\times\RC U;\Sphere V\times\Sphere U]
\end{equation*}
onto the closure of $E(W)$. Hence the identity of $W$ extends to a diffeomorphism
$\mathcal D_1\cong\mathcal D_2$, and both are the doubly-relative compactification
$\relRC{U,V}{U\times V}$. \textup{(}In $\mathcal D_1$ the centres $\Sphere U$, $\Sphere V\subset\Sphere W$
are disjoint, so the order does not matter.\textup{)}
\end{proposition}
\begin{proof}
Write $\rho_U=\ang u^{-1}$, $\rho_V=\ang v^{-1}$ and note $\ang w^2=\ang u^2+\ang v^2-1$. If
$\bar E:\mathcal D\to\R^N$ is a continuous extension of $E$ to a compact space $\mathcal D$ in
which $W$ is dense, then $\bar E(\mathcal D)$ is compact, contains $E(W)$ densely, and hence
equals $\overline{E(W)}$. So it suffices to show that $E$ extends to injective immersions of
$\mathcal D_1$ and $\mathcal D_2$; an injective immersion of a compact manifold with
corners is a diffeomorphism onto its image.

\emph{$\mathcal D_2$.} The components $\rho_U$, $u/\ang u$, $\rho_V$, $v/\ang v$ of $E$ are
smooth on $\RC V\times\RC U$, and together they embed it. The others are smooth functions of
these and of $\rho_V/\sqrt Q$, $\rho_U/\sqrt Q$, where
$Q=\rho_U^2+\rho_V^2-\rho_U^2\rho_V^2$: multiplying numerator and denominator by
$\rho_U\rho_V$ gives $\ang u/\ang w=\rho_V/\sqrt Q$, $\ang v/\ang w=\rho_U/\sqrt Q$, and then
$v/\ang w=(v/\ang v)\rho_U/\sqrt Q$, $u/\ang w=(u/\ang u)\rho_V/\sqrt Q$. Since $Q>0$ off the corner, $E$ extends to an
embedding of $\mathcal D_2\setminus\mathrm{ff}$. Near the front face use polar coordinates
$\rho_V=R\cos\psi$, $\rho_U=R\sin\psi$: then
$Q=R^2(1-R^2\cos^2\psi\sin^2\psi)$, so the two functions are smooth, and at $R=0$ the
components $\ang u/\ang w=\cos\psi$, $\ang v/\ang w=\sin\psi$ have differentials
$-\sin\psi\,d\psi$, $\cos\psi\,d\psi$ while $d\rho_V=\cos\psi\,dR$, $d\rho_U=\sin\psi\,dR$; with the
directions $v/\ang v$, $u/\ang u$ this shows $\bar E$ is an immersion. It is injective:
$\rho_U,u/\ang u,\rho_V,v/\ang v$ determine the image in $\RC V\times\RC U$, and on a fibre of
the front face $\psi$ is determined by $\ang u/\ang w$, $\ang v/\ang w$.

\emph{$\mathcal D_1$.} By \cref{blowup:thm:relative},
$[\RC W;\Sphere V]\cong\RC U\times\RC{V}$, with interior coordinates $(u,q)$, $q=v/\ang u$.
The components $\rho_U$, $u/\ang u$, $\ang u/\ang w=1/\ang q$ and $v/\ang w=q/\ang q$ of $E$
are smooth on it and embed it; the remaining components $(v,u)/\ang w$ are smooth on
$\RC W$, while $\rho_V$, $v/\ang v$, $\ang v/\ang w$ are smooth on $[\RC W;\Sphere U]$ by
symmetry. Since $\Sphere U\cap\Sphere V=\emptyset$, $\mathcal D_1$ maps smoothly to $[\RC
W;\Sphere V]$, to $[\RC W;\Sphere U]$ and to $\RC W$, so $\bar E$ is smooth on $\mathcal D_1$, and
it is an injective immersion on $\mathcal D_1\setminus\mathrm{ff}_U\cong[\RC W;\Sphere V]\setminus\Sphere U$
and, symmetrically, on $\mathcal D_1\setminus\mathrm{ff}_V$. These two open sets cover $\mathcal
D_1$. If $\bar E(p)=\bar E(p')$, the functions $1/\ang w=\rho_U\cdot\ang u/\ang w$ and
$w/\ang w$, which are determined by $\bar E$, show that $\beta(p)=\beta(p')$ in $\RC W$; then $p,p'$ lie in the same one of the two open sets (both in
$\mathrm{ff}_U$, both in $\mathrm{ff}_V$, or neither), so $p=p'$.
\end{proof}

%======================================================================
\section{Projective blow-up and blow-up of hypersurfaces}\label{blowup:sec:projective}

In \emph{projective} blow-up one uses two-sided polar coordinates, identifying $\omega$
with $-\omega$. For the origin in $\R^n$ this gives the tautological line bundle
\begin{equation*}
  L=\{(z,\ell)\in\R^n\times\R\mathbb P^{n-1};\ z\in\ell\},\qquad\beta_{\mathbb P}(z,\ell)=z.
\end{equation*}

\begin{proposition}\label{blowup:prop:projective}
$L$ is a smooth line bundle over $\R\mathbb P^{n-1}$, $\beta_{\mathbb P}$ is smooth and maps
$L\setminus0_L$ diffeomorphically onto $\R^n\setminus0$, and on $\{\ell_j\ne0\}$ the functions
$(z_j,t_k=z_k/z_j)$, $z_j\in\R$, are coordinates. The map
$[0,\infty)\times\Sphere^{n-1}\ni(r,\omega)\mapsto(r\omega,[\omega])\in L$ induces a
diffeomorphism $[\R^n;\{0\}]\cong[L;0_L]$; thus radial blow-up is projective blow-up
followed by the blow-up of the hypersurface $0_L\cong\R\mathbb P^{n-1}$.
\end{proposition}
\begin{proof}
Over $\{\ell_j\ne0\}$ the section $\ell\mapsto(t_k)$ with $t_j=1$ trivializes $L$, and
$z=z_j(t,1)$, so $(z_j,t)\in\R\times\R^{n-1}$ are coordinates on $L$, with $\beta_{\mathbb P}$
smooth; for $z\ne0$ the inverse $z\mapsto(z,[z])$ is smooth. By
\cref{blowup:prop:bundle}, $[L;0_L]\cong[0,\infty)\times SL$, using the metric from $\R^n$,
and $SL=\{(\omega,[\omega]);\ \omega\in\Sphere^{n-1}\}\cong\Sphere^{n-1}$. In these terms the
given map is $(r,\omega)\mapsto(r,(\omega,[\omega]))$.
\end{proof}

The projective blow-up has no boundary, but it is not orientable for even $n$, and the
functions we need are usually smooth only from each side of $0_L$: $\abs z=\abs{z_j}(1+\abs
t^2)^{1/2}$ is not smooth across $z_j=0$. Blowing up the one-sided hypersurface $0_L$
separates the two sides, which is why radial blow-up is preferred. For a two-sided
hypersurface blow-up simply cuts the space open.

\begin{proposition}\label{blowup:prop:hypersurface}
Let $H\subset X$ be a closed interior p-submanifold of codimension one in a manifold with
corners. Then the front face of $[X;H]$ is the double cover $SNH$ of $H$. If
$H=\{f=0\}$ with $f\in\CI(X;\R)$ and $df\ne0$ on $H$, then $[X;H]$ is canonically the
disjoint union $X_+\sqcup X_-$ of the manifolds with corners $X_\pm=\{\pm f\ge0\}$, the
diffeomorphism being the identity on $\{f\ne0\}$.
\end{proposition}
\begin{proof}
The fibre $\Sphere^{0,1}=\{\pm1\}$ of the front face has two points. In an adapted chart
$H=\{x_1=0\}$ with $x_1$ an interior coordinate, and $f=x_1g$ with $g\ne0$ near $H$ since
$df\ne0$ there; so $(f,y)$ is again an adapted chart. In it, polar coordinates are
$(r,\omega)=(\abs f,\sgn f)$, and the map $X_\pm\ni(f,y)\mapsto(\pm f,\pm1,y)$ is a
diffeomorphism onto $\{\omega=\pm1\}$.
\end{proof}

%======================================================================
\section{b-maps, b-submersions and b-fibrations}\label{blowup:sec:bfib}

\subsection{Fibrations of manifolds with corners}
A smooth map $f:X\to X'$ of manifolds with corners is a \emph{fibration} if each
$p'\in f(X)$ has an open neighbourhood $U$ with a diffeomorphism $f^{-1}(U)\cong U\times F$
under which $f$ becomes the projection. For compact manifolds without boundary, a
surjective submersion is a fibration (Ehresmann's theorem). For manifolds with corners
this fails: the inclusion $[0,1]\hookrightarrow\R$ has surjective differential, and the
projection $(x,y)\mapsto x$ of the closed half-disc
$\{x^2+y^2\le1,\ y\ge0\}$ onto $[-1,1]$ is surjective, with surjective differential
everywhere, but its fibres are intervals over $(-1,1)$ and points over $\pm1$
(\cref{blowup:fig:fibrations}, \cref{blowup:ex:half-disc}). The correct condition involves
the boundary faces.

\begin{proposition}\label{blowup:prop:fibration}
Let $X$ be a compact manifold with corners, $X'$ a manifold without boundary and $f:X\to
X'$ smooth. If the restriction of $f$ to every boundary face of $X$ \textup{(}including
$X$ itself\textup{)} is a submersion, then $f$ is a fibration over its image, which is open
and closed; its fibres are compact manifolds with corners. Conversely, a fibration over
$X'$ has this property.
\end{proposition}
\begin{proof}
Fix $p'\in f(X)$ and coordinates $y'$ near $p'$ on a ball $B$, $y'(p')=0$. We construct
vector fields $V_i$ on $f^{-1}(B)$, tangent to every boundary hypersurface, with
$f_*V_i=\pa_{y'_i}$. Near $p\in f^{-1}(B)$ of codimension $l$ take a chart $(x,y)\in\R^{n,l}$
centred at $p$; the boundary face through $p$ is $\{x=0\}$, and the hypothesis says that
$D_yf(0,0)$ is surjective, hence $D_yf(x,y)$ is surjective nearby. Put
$W_i=\sum_jc_{ij}\pa_{y_j}$ with $c=(D_yf)^T(D_yf\,(D_yf)^T)^{-1}$ applied to $e_i$. These
satisfy $f_*W_i=\pa_{y'_i}$ and have no $\pa_x$ components, so they lie in $\Vb$. Both
conditions are preserved by combinations with a partition of unity, so we get $V_i$ on
$f^{-1}(B)$. A vector field in $\Vb$ on a compact manifold with corners has a flow defined
for all times that preserves $X$: in local coordinates the $x_i$-component of an integral
curve satisfies $\dot x_i=a_ix_i$, so $x_i(t)=x_i(0)\exp\int_0^ta_i$ keeps its sign.
(Here we apply this to $\chi V_i$, with $\chi=1$ on $f^{-1}(B')$, $B'\Subset B$, and
$\chi\in\CI(X)$ supported in $f^{-1}(B)$.) Let $B''$ be a ball around $0$ of radius less than half the radius of $B'$. For
$y'\in B''$ let $\Phi_{y'}$ be the time-one flow of $\sum y'_i\chi V_i$. Where $\chi=1$ this
field is $f$-related to the constant field $\sum y'_i\pa_{y'_i}$, so for $f(q)\in B''$ the
image under $f$ of the integral curve moves along a straight segment inside $B'$, and
$f(\Phi_{y'}(q))=f(q)+y'$. Then, for $B'''\subset B''$ a smaller ball,
\begin{equation*}
  B'''\times f^{-1}(p')\ni(y',q)\longmapsto\Phi_{y'}(q)\in f^{-1}(B'''),\qquad
  q'\longmapsto\big(f(q'),\Phi_{-f(q')}(q')\big),
\end{equation*}
are inverse smooth maps, and $f$ becomes the projection. The fibre $f^{-1}(p')$ is a
manifold with corners: in the chart above it is $\{f(x,y)=p'\}$, and since $D_yf$ is
surjective the implicit function theorem (with $x$ as parameters) solves for some of the
$y$. This also shows $f(X)$ is open; it is closed since $X$ is compact. Conversely, in a
local product $U\times F$ the boundary faces are $U\times G$, $G$ a face of $F$, and the
projection restricted to them is a submersion.
\end{proof}

\begin{figure}[ht]
\centering
\begin{tikzpicture}[>=Stealth,scale=1.0]
  \begin{scope}
    \fill[gray!12] (-1.5,0) arc (180:0:1.5) -- cycle;
    \draw[very thick] (-1.5,0) arc (180:0:1.5) -- cycle;
    \foreach \x in {-1.2,-0.6,0,0.6,1.2}{
      \pgfmathsetmacro{\h}{sqrt(2.25-\x*\x)}
      \draw[gray] (\x,0)--(\x,\h);}
    \draw[->] (0,-0.4)--(0,-0.9);
    \draw[thick] (-1.5,-1.2)--(1.5,-1.2);
    \fill (-1.5,-1.2) circle (1.5pt) (1.5,-1.2) circle (1.5pt);
    \node[below,font=\footnotesize] at (-1.5,-1.25) {$-1$};
    \node[below,font=\footnotesize] at (1.5,-1.25) {$1$};
    \node[font=\small,align=center] at (0,-2.3) {(a) surjective, $df$ surjective,\\ not a fibration};
  \end{scope}
  \begin{scope}[xshift=5.5cm]
    \fill[gray!12] (0,0) rectangle (2.4,2.4);
    \fill[white] (0,0) -- (0.8,0) arc (0:90:0.8) -- cycle;
    \draw[very thick] (0,2.4)--(0,0.8) arc (90:0:0.8) -- (2.4,0) -- (2.4,2.4) -- cycle;
    \draw[line width=2.2pt] (0,2.4)--(0,0.8) arc (90:0:0.8);
    \foreach \x in {0.3,0.6}{
      \pgfmathsetmacro{\h}{sqrt(0.64-\x*\x)}
      \draw[gray] (\x,\h)--(\x,2.4);}
    \foreach \x in {1.1,1.6,2.0} \draw[gray] (\x,0)--(\x,2.4);
    \draw[->] (1.2,-0.4)--(1.2,-0.9);
    \draw[thick] (0,-1.2)--(2.4,-1.2);
    \fill (0,-1.2) circle (1.5pt);
    \node[below,font=\footnotesize] at (0,-1.25) {$0$};
    \node[font=\small,align=center] at (1.2,-2.3) {(b) $[[0,1]^2;\{0\}]\to[0,1]$:\\ a b-fibration};
  \end{scope}
\end{tikzpicture}
\caption{(a) The projection of the half-disc onto $[-1,1]$: the fibres (gray) degenerate
to points at the ends. (b) The blow-down of the corner of $[0,1]^2$ followed by
$(x,y)\mapsto x$: the fibres over $x>0$ are segments, while the fibre over $0$ (heavy) is
the union of the front face and the lift of $\{x=0\}$, two boundary hypersurfaces
meeting at a corner.}
\label{blowup:fig:fibrations}
\end{figure}

When $X'$ has boundary, maps such as blow-downs, which are not fibrations, must also be
allowed. The right notion uses the b-tangent bundle.

\subsection{The b-tangent bundle and the b-differential}
By \cref{blowup:lem:Vb}, $\Vb(X)$ is locally free with basis $x_i\pa_{x_i}$, $\pa_{y_j}$.
Hence there is a vector bundle $\bT X$, the \emph{b-tangent bundle}, with
$\bT_pX=\Vb(X)/\mathcal I_p\Vb(X)$ ($\mathcal I_p$ the functions vanishing at $p$) and
$\CI(X;\bT X)=\Vb(X)$. For $v=V_p\in\bT_pX$ and a function $h$ which is smooth or of the
form $\log(a\prod\rho_H^{e_H})$ with $a>0$ smooth, set $v(h)=(Vh)(p)$; this is well defined
because $V\log\rho_H=V\rho_H/\rho_H$ is smooth for $V\in\Vb(X)$, and it vanishes if $V\in\mathcal
I_p\Vb$.

\begin{lemma}\label{blowup:lem:bdiff}
Let $f:X\to X'$ be an interior b-map and $p\in X$, $p'=f(p)$. There is a unique linear map
$f_*:\bT_pX\to\bT_{p'}X'$, the \emph{b-differential}, such that $(f_*v)(g)=v(f^*g)$ for all
$g$ smooth near $p'$ and $(f_*v)(\log\rho_{H'})=v(\log f^*\rho_{H'})$ for all $H'\ni p'$.
For composable interior b-maps, $(g\circ f)_*=g_*f_*$.
\end{lemma}
\begin{proof}
In a chart $(x',y')$ at $p'$, an element $\bar v=\sum\alpha_ix'_i\pa_{x'_i}+\sum\beta_j\pa_{y'_j}$
of $\bT_{p'}X'$ is determined by $\bar v(\log x'_i)=\alpha_i$ and $\bar v(y'_j)=\beta_j$; this
gives uniqueness. Define $f_*v$ by $\alpha_i=v(\log f^*x'_i)$, $\beta_j=v(f^*y'_j)$, which
makes sense by \eqref{blowup:eq:bmap}. For smooth $g$, the chain rule on the dense set
$X^\circ$ gives $V(f^*g)=\sum f^*(\pa_{x'_i}g)\,f^*x'_i\,V(\log f^*x'_i)+\sum
f^*(\pa_{y'_j}g)V(f^*y'_j)$, whose value at $p$ is $(f_*v)(g)$. For a defining function
$\rho_{H'}=cx'_i$ we have $\log\rho_{H'}=\log c+\log x'_i$, and both sides are additive. The
composition rule follows from uniqueness, since $v(f^*g^*h)=(f_*v)(g^*h)=(g_*f_*v)(h)$ for
both kinds of $h$.
\end{proof}

\begin{definition}\label{blowup:def:bfib}
An interior b-map $f:X\to X'$ is a \emph{b-submersion} if $f_*:\bT_pX\to\bT_{f(p)}X'$ is
surjective for all $p$. It is \emph{b-normal} if for each $H\in\mathcal M_1(X)$ there is at
most one $H'\in\mathcal M_1(X')$ with $e(H,H')\ne0$, i.e.\ no boundary hypersurface is mapped
into a corner of codimension two or more. A \emph{b-fibration} is a b-normal
b-submersion.
\end{definition}

A fibration between manifolds with corners which is an interior b-map is a b-fibration,
by the local product form. Composites of b-submersions are b-submersions by
\cref{blowup:lem:bdiff}. The blow-down map of an interior p-submanifold is not a
b-submersion: for $[\R^2;\{0\}]$, $\bT=T$ downstairs and $\beta_*(r\pa_r)=x\pa_x+y\pa_y$ and
$\beta_*(\pa_\omega)=x\pa_y-y\pa_x$ both vanish at $0$.

\begin{proposition}\label{blowup:prop:bd-bsub}
For a boundary face $F$ of a manifold with corners $X$, the blow-down map
$\beta:[X;F]\to X$ is a b-submersion; in fact $\beta_*$ is an isomorphism on each fibre of
$\bT$. If $F$ has codimension at least two, $\beta$ is not b-normal.
\end{proposition}
\begin{proof}
Away from the front face $\beta$ is a local diffeomorphism. Near the front face use an
adapted chart with $F=\{x=0\}$, $x\in[0,\infty)^a$, and projective coordinates
$s=x_j$, $t_m=x_m/x_j$ ($m\ne j$), all of them boundary coordinates. Then
$\log\beta^*x_j=\log s$ and $\log\beta^*x_m=\log s+\log t_m$, so
\begin{equation*}
  \beta_*(s\pa_s)=\sum_{i=1}^ax_i\pa_{x_i},\qquad\beta_*(t_m\pa_{t_m})=x_m\pa_{x_m},\qquad
  \beta_*(\pa_{y})=\pa_y,
\end{equation*}
and these images form a basis of $\bT X$. If $a\ge2$, \cref{blowup:prop:props}(ii) gives
$e(\mathrm{ff},H_i)=1$ for the $a\ge2$ hypersurfaces $H_i=\{x_i=0\}$ containing $F$.
\end{proof}

\begin{example}\label{blowup:exa:bfib}
The composite $f=\pi_1\circ\beta:[[0,1]^2;\{(0,0)\}]\to[0,1]$ of the blow-down of the corner
and the projection $\pi_1(x,y)=x$ is a b-fibration which is not a fibration
(\cref{blowup:fig:fibrations}(b)). It is a b-submersion as a composite of b-submersions
($\pi_1$ maps $x\pa_x\mapsto x\pa_x$ at $x=0$ and $\pa_x\mapsto\pa_x$ elsewhere). It is
b-normal because $[0,1]$ has no corners of codimension two. In the projective coordinates
$s=y$, $t=x/y$ near the corner $q$ of the front face and the lift of $\{x=0\}$,
$f^*x=st$, whose differential vanishes at $q$. If $f$ were a fibration near $q$, then in a
product decomposition $[0,\delta)\times F$ the function $f^*x$ would be the first coordinate,
with non-vanishing differential. Note that the fibre over $0$ is the union of two boundary
hypersurfaces.
\end{example}

\begin{discussion}
Why are b-fibrations the right replacement for fibrations? The b-differential sees the
maps $x\mapsto x^e$ in boundary variables as ``linear'' ($\log x\mapsto e\log x$), so a
b-submersion is a submersion in the logarithmic coordinates $\log x_i$, which is what
matters for integration along the fibres of functions with expansions in powers and logarithms
of boundary defining functions. b-normality prevents a whole boundary hypersurface from
being squeezed into a corner, where the expansions of different faces would collide.
Melrose's push-forward theorem (see the Notes) makes this precise.
\end{discussion}

%======================================================================
\section{A dictionary of blown-up spaces in this book}\label{blowup:sec:dictionary}

We list the spaces of the book which are blow-ups, with the results above which justify
the descriptions given where they are used.

\begin{enumerate}[label=(\arabic*),leftmargin=*]
\item \emph{Radial compactification} (\cref{ch:distributions}). In the inverted coordinate
$w=4z/\abs z^2$ near $\infty\in\oRC{\R^n}$, the collar coordinates of
\cref{distributions:lem:collar} are $s=\abs w/4$, $\omega=w/\abs w$: polar coordinates in
$w$. So $\RC{\R^n}=[\oRC{\R^n};\{\infty\}]$, and this description is independent of
coordinates by \cref{blowup:thm:invariance}; this answers \cref{distributions:ex:polar}.
\item \emph{Blow-up of the origin} $[W;\{0\}]$ (\cref{conormal:lem:blowup},
\cref{conormal:ex:blowup}): \cref{blowup:prop:bundle}, with \cref{blowup:exa:homogeneous}.
\item \emph{Relative compactifications.} $\relRC VW=[\RC W;\Sphere V]$ with $H_V$ the front
face (\cref{distributions:rem:blowup}, \cref{product:prop:blow-down},
\cref{product:ex:blow-up}): \cref{blowup:thm:relative}. The doubly-relative
compactification is $[\RC{V\times U};\Sphere U,\Sphere V]=[\RC V\times\RC U;\Sphere V\times\Sphere U]$
(\cref{distributions:rem:blowup}, \cref{distributions:ex:doubly}):
\cref{blowup:prop:doubly}.
\item \emph{Product-type symbols} (\cref{ch:product}). Fibrewise, for a subbundle
$V=\phi^*T^*B$ of $W=T^*M$, the relative compactification is
$[\RC{T^*M};\phi^*S^*B]$, with main face $H_W=[S^*M;\phi^*S^*B]$ by
\cref{blowup:thm:relative,blowup:prop:lift}, fibred over $S^*(M/B)$ with fibres radially
compactified affine spaces by \cref{blowup:lem:sphere-fibration}.
\item \emph{Families symbols} (\cref{ch:families}). The symbol space
$[S^*(\tilde M/B);\Sigma]$, $\Sigma=S(\tilde\phi^*T^*(M/B))$, is the blow-up of an interior
submanifold of a manifold without boundary; the map
$\gamma:[S^*(\tilde M/B);\Sigma]\to S^*(\tilde M/M)$ in \ref{families:P-fibre} is the map of
\cref{blowup:lem:sphere-fibration} for $W=T^*(\tilde M/B)$, $V=\tilde\phi^*T^*(M/B)$,
$W/V=T^*(\tilde M/M)$.
\item \emph{The circle bundle in \cref{ch:eta}.} In \ref{eta:P1}, $S^*M\times\Sphere=
\{\tau=0\}$ is a two-sided hypersurface of $S^*(M\times\Sphere)$: on the unit cosphere
$\abs\xi^2+\tau^2=1$, $d\tau\ne0$ where $\tau=0$. By \cref{blowup:prop:hypersurface} the
blow-up is the disjoint union of the closed half-sphere bundles $S^*_\pm=\{\pm\tau\ge0\}$.
\item \emph{Semiclassical and adiabatic double spaces} (\cref{ch:semiclassical}):
$[[0,1]\times X^2;\{0\}\times\Diag]$ and $[[0,1]\times Y^2;\{0\}\times Y^2_\phi]$, described in
\cref{blowup:prop:sc-kernels}, \cref{blowup:cor:sc-manifold} and
\cref{blowup:rem:adiabatic}; these justify \cref{semiclassical:rem:double-space} and
\cref{semiclassical:prop:smoothing-X}(iii).
\end{enumerate}

\begin{example}\label{blowup:exa:b-double}
Let $X$ be a manifold with boundary. The diagonal of $X^2$ is not a p-submanifold
(\cref{blowup:ex:diag-not-p}), but it lifts to one in the \emph{b-double space}
$X^2_{\mathrm b}=[X^2;(\pa X)^2]$ (for $\pa X$ connected; in general one blows up the
products $H_i\times H_i$ of the boundary components). In coordinates $(x,y)$, $(x',y')$
on the two factors, $(\pa X)^2=\{x=x'=0\}$ is a corner of codimension two. In the
projective coordinates $s=x$, $t=x'/x$ the diagonal minus $(\pa X)^2$ is $\{t=1,\ y=y'\}$
with $s>0$, so its closure $\Diag_{\mathrm b}=\{t=1,\ y=y'\}$ is an interior p-submanifold
meeting the front face transversally (\cref{blowup:fig:corner}). The b-vector fields of the
left factor lift to $x\pa_x=s\pa_s-t\pa_t$ and $\pa_{y_j}$, and since
$(s\pa_s-t\pa_t)(t-1)=-1$ and $\pa_{y_j}(y_k-y'_k)=\delta_{jk}$ on $\Diag_{\mathrm b}$, they are
transversal to $\Diag_{\mathrm b}$: they span its normal bundle at every point. This is the
starting point of the b-calculus \cite{Melrose-APS}, in which the kernels of
b-pseudodifferential operators are conormal to $\Diag_{\mathrm b}$ in $X^2_{\mathrm b}$.
\end{example}

\begin{remark}\label{blowup:rem:other-calculi}
Other calculi of pseudodifferential operators on manifolds with boundary arise in the
same way, by resolving the diagonal together with a Lie algebra of vector fields: the
$0$-calculus of Mazzeo and Melrose (vector fields vanishing at $\pa X$, double space
$[X^2;\pa\Diag]$), the scattering calculus (double space $[X^2_{\mathrm b};\pa\Diag_{\mathrm b}]$),
and the cusp and fibred-boundary calculi. We do not use them; see \cite{Melrose-APS} for
the b-calculus and the Notes for further references.
\end{remark}

%======================================================================
\section{Exercises}\label{blowup:sec:exercises}

\begin{exercise}\label{blowup:ex:projective}
In the projective coordinates $s=z_1$, $t_k=z_k/z_1$ ($k\ge2$) on
$\{\omega_1>0\}\subset[\R^n;\{0\}]$: (a) verify \eqref{blowup:eq:proj-lifts}; (b) show that the
Euler field lifts to $s\pa_s$ and the rotation $z_1\pa_{z_2}-z_2\pa_{z_1}$ to
$(1+t_2^2)\pa_{t_2}-t_2s\pa_s+t_2\sum_{r\ge3}t_r\pa_{t_r}$; (c) show that $\pa_{z_1}$ lifts to
$s^{-1}$ times a smooth vector field, but not to a smooth one; (d) compute
$\beta^*\abs{dz}$ in these coordinates.
\end{exercise}

\begin{exercise}\label{blowup:ex:cone}
Let $C=\{(t,x,y)\in\R^3;\ t^2=x^2+y^2,\ t\ge0\}$. Show that the closure of
$\beta^{-1}(C\setminus0)$ in $[\R^3;\{0\}]$ is an interior p-submanifold diffeomorphic to
$[0,\infty)\times\Sphere^1$, meeting the front face transversally in a circle, and that
$\beta$ restricts to the polar map $(r,\theta)\mapsto r(1,\theta)/\sqrt2$ on it.
\end{exercise}

\begin{exercise}\label{blowup:ex:morse}
Let $u\in\CI(\R^n)$ have a non-degenerate critical point at $0$, so that
$u(z)=u(0)+z^TH(z)z$ with $H$ smooth, symmetric, and $H(0)$ invertible. Show that on
$[\R^n;\{0\}]$, $\beta^*u=u(0)+r^2U$ with $U(r,\omega)=\omega^TH(r\omega)\omega$, and that
$dU\ne0$ near $\{r=0,\ U=0\}$. Conclude that near the front face the lift of the level set
$\{u=u(0)\}$ is the union of the front face and the hypersurface $\{U=0\}$, which meet
transversally.
\end{exercise}

\begin{exercise}\label{blowup:ex:diag-not-p}
Let $X$ be a manifold with boundary of dimension $n\ge1$ with $\pa X\ne\emptyset$. Show that
$\Diag\subset X^2$ is not a p-submanifold. (Count the dimensions of $\Diag\cap\{x=0\}$,
$\Diag\cap\{x'=0\}$ and $\Diag\cap\{x=x'=0\}$.)
\end{exercise}

\begin{exercise}\label{blowup:ex:noncommute}
In $\R^3$ let $Y=\{z_1=z_2=0\}$ and $Z=\{z_2=z_3=0\}$. Show that $Y$ and $Z$ are joint
p-submanifolds, neither transversal nor nested, and that $f=z_2/\abs{(z_1,z_2)}$ extends
continuously to $[\R^3;Y;Z]$ but not to $[\R^3;Z;Y]$. Conclude that the identity of
$\R^3\setminus(Y\cup Z)$ does not extend to a homeomorphism $[\R^3;Z;Y]\to[\R^3;Y;Z]$.
\end{exercise}

\begin{exercise}\label{blowup:ex:corner-exponents}
Let $\beta:[\R^{2,2};\{0\}]\to\R^{2,2}=[0,\infty)^2$ be the blow-down of the corner. Compute
the exponent matrix of $\beta$ and verify directly that $\beta$ is a b-submersion which is
not b-normal, and that $\pi_1\circ\beta$ is b-normal. More generally, show that the blow-down of
a boundary face of codimension $a$ has $e(\mathrm{ff},H)=1$ for the $a$ hypersurfaces
$H\supset F$.
\end{exercise}

\begin{exercise}\label{blowup:ex:half-disc}
Let $D_+=\{(x,y);\ x^2+y^2\le1,\ y\ge0\}$ and $f(x,y)=x$. Show that $D_+$ is a manifold with
corners, that $f:D_+\to[-1,1]$ is surjective with surjective differential, and that $f$ is
not a fibration. Which hypothesis of \cref{blowup:prop:fibration} fails for
$f:D_+\to\R$? Is $f:D_+\to[-1,1]$ a b-map?
\end{exercise}

\begin{exercise}\label{blowup:ex:octagon}
Let $\dim U=\dim V=1$. Show that $\relRC{U,V}{U\times V}$ is an octagon
(\cref{blowup:fig:octagon}): it has eight boundary hypersurfaces, each an interval.
Identify them in both descriptions of \cref{blowup:prop:doubly}. Show that if $\dim U,\dim
V\ge2$ there are exactly three boundary hypersurfaces, and identify them.
\end{exercise}

\section*{Notes}
This appendix is based on Lectures~1 and~2 of Melrose's MSRI lectures \emph{Real blow
ups} (2008); the remaining sections of those notes were not written and are omitted, apart
from the brief \cref{blowup:exa:b-double,blowup:rem:other-calculi}. Several statements
have been corrected or completed. The coordinate invariance of polar coordinates is proved
by an explicit Taylor-formula lift (\cref{blowup:thm:invariance}) instead of the
homotopy argument sketched originally, and it is proved with parameters and at corners.
The lifts of linear vector fields in projective coordinates
\eqref{blowup:eq:proj-lifts} correct the formula of the source. The front face is
$S^+NY$ with $NY=T_YX/TY$ (not $TY$), the blow-up of the zero section is $[E;0_E]$, and
the reduction of the general case to the vector bundle case uses the tubular (not collar)
neighbourhood theorem. The projective blow-up of a point is the tautological line bundle
over $\R\mathbb P^{n-1}$. The commutation results for nested and transversal centres are
proved by identifying both iterated blow-ups with an explicit model, and a counterexample
shows that commutation fails in general. The source asserts that a surjective map with
surjective differential between manifolds with corners is a fibration; the half-disc shows
this is false, and \cref{blowup:prop:fibration} gives a correct criterion. The
identifications of the relative and doubly-relative compactifications, the semiclassical
criterion and the sphere fibration lemma were added for use in the book. Real blow-up is
developed systematically in Melrose's unpublished book \emph{Differential analysis on
manifolds with corners}; b-fibrations and the push-forward theorem for polyhomogeneous
functions are in R.~B. Melrose, \emph{Calculus of conormal distributions on manifolds with
corners}, Int.\ Math.\ Res.\ Notices 1992; the b-calculus is in \cite{Melrose-APS}; the
$0$-calculus is in R.~Mazzeo and R.~B. Melrose, J.\ Funct.\ Anal.\ 75 (1987).
